% Môn: Vật lý
% Lớp: 11
% Chương: Chương I. Dao động
% Bài: L11C1 Ôn Tập chương dao động
% Số câu: 137
% App đọc file này trực tiếp — sửa rồi Commit trên GitHub, bấm Đồng bộ GitHub .tex

% L11C1 Ôn Tập chương dao động

% ===== Câu 1 =====
% ID: L11C1B1-29-DS
% Mức: TH
\dangbt{quãng đường}
\begin{ex}
% ID: L11C1B1-29-DS
% Mức: TH
L11C1 Ôn Tập chương dao động – quãng đường: Vật có biên độ $6\,\text{cm}$, chu kì $1\,\text{s}$.
\choiceTF
{\True Một chu kì đi $24\,\text{cm}$.}
{\True Nửa chu kì đi $12\,\text{cm}$.}
{\True Sau một chu kì độ dịch chuyển bằng 0.}
{\True Tốc độ trung bình là $24\,\text{cm}$ /s.}
\end{ex}

% ===== Câu 2 =====
% ID: L11C1B1-30-DS
% Mức: TH
\dangbt{Chưa có dạng}
\begin{ex}
% ID: L11C1B1-30-DS
% Mức: TH
Một vật dao động điều hoà với chu kì T. Tại thời điểm ban đầu, vật đi qua vị trí cân bằng. Tính $T$ tỉ số giữa động năng và thế năng của vật vào thời điểm 12. $\nguon{ly11 kntt.pdf}$
\choiceTF

\loigiai{
Theo giản đồ Hình I.2 $G$, ta thấy khi t = $\dfrac{T}{12}$ thì x = $\pm \dfrac{A}{2}$

[Một vật dao động điều hoà với chu kì T. Tại thời điểm ban đầu]

$W_{t} = \dfrac{1}{2}kx^{2} = \dfrac{1}{2}k\left( \dfrac{A}{2} \right)^{2} = \dfrac{1}{4}k\dfrac{A^{2}}{2} = \dfrac{1}{4}W\left. \text{W}_{\text{d}} = \text{W} - \text{W}_{\text{t}} = \text{W} - \dfrac{1}{4}\text{W} = \dfrac{3}{4}\text{W}\Rightarrow\dfrac{\text{W}_{\text{d}}}{\text{W}_{\text{t}}} = 3 \right.$.
}
\end{ex}

% ===== Câu 3 =====
% ID: L11C1B1-31-TLN
% Mức: TH
\dangbt{quãng đường}
\begin{ex}
% ID: L11C1B1-31-TLN
% Mức: TH
L11C1 Ôn Tập chương dao động – quãng đường: Vật có biên độ $9\,\text{cm}$. Quãng đường trong 3 chu kì theo cm bằng bao nhiêu?
\shortans{108}
\end{ex}

% ===== Câu 4 =====
% ID: L11C1B1-32-TN
% Mức: NB
\dangbt{Chưa có dạng}
\begin{ex}
% ID: L11C1B1-32-TN
% Mức: NB
Một vật đang dao động điều hoà dưới tác dụng của một lực đàn hồi. Chọn câu đúng. $\nguon{ly11 kntt.pdf}$
\choice
{Khi vật đi qua vị trí cân bằng thì gia tốc đạt giá trị cực đại.}
{Khi vật ở vị trí biên thì lực đổi chiều.}
{\True Khi vật đi từ vị trí cân bằng đến vị trí biên thì gia tốc ngược chiều với vận tốc.}
{Khi vật đi từ vị trí biên đến vị trí cân bằng thì độ lớn của gia tốc tăng dần.}
\loigiai{
Chọn C.

Khi vật đi từ vị trí cân bằng đến vị trí biên thì vận tốc hướng về vị trí biên còn gia tốc vẫn luôn hướng về vị trí cân bằng.
}
\end{ex}

% ===== Câu 5 =====
% ID: L11C1B1-33-DS
% Mức: VD
\dangbt{quãng đường}
\begin{ex}
% ID: L11C1B1-33-DS
% Mức: VD
L11C1 Ôn Tập chương dao động – quãng đường: Vật có biên độ $7\,\text{cm}$, chu kì $2\,\text{s}$.
\choiceTF
{\True Một chu kì đi $28\,\text{cm}$.}
{\True Nửa chu kì đi $14\,\text{cm}$.}
{\True Sau một chu kì độ dịch chuyển bằng 0.}
{\True Tốc độ trung bình là $14\,\text{cm}$ /s.}
\end{ex}

% ===== Câu 6 =====
% ID: L11C1B1-34-DS
% Mức: TH
\dangbt{Chưa có dạng}
\begin{ex}
% ID: L11C1B1-34-DS
% Mức: TH
Một vật dao động điều hoà dọc theo trục Ox nằm ngang, gốc $O$ và mốc thế năng ở vị trí cân bằng. Cứ sau $0,5\,\mathrm{s}$ thì động năng lại bằng thế năng và vật đi được đoạn đường dài nhất trong thời gian $0,5\,\mathrm{s}$ là $4\sqrt{2} \mathrm{cm}$. Chọn $t = 0 là lúc vật qua vị trí cân bằng theo chiều dương$. Viết phương trình dao động của vật. $\nguon{ly11 kntt.pdf}$
\choiceTF

\loigiai{
Khi $. W_{t} = W_{d}\Rightarrow W_{t} = \dfrac{W}{2} ..\Leftrightarrow\dfrac{kx^{2}}{2}=\dfrac{1}{2}\dfrac{kA^{2}}{2}\Rightarrowx =$ \pm\dfrac{A\sqrt{2}}{2}\right $.

[Một vật dao động điều hoà dọc theo trục Ox nằm ngang]

Từ giản đồ ở Hình I.2Ga, ta có:$\left$.$\dfrac{T}{4}$= 0,5$\RightarrowT$= 4.0,5 = 2$\text{s}\right$.$\left. \Rightarrow\omega = \dfrac{2\pi}{T} = \pi\text$(rad/s)$\right.$Từ giản đồ Hình I.2Gb, ta có:$\text{S}_{\max} = A\sqrt{2} = 4\sqrt{2}\text{cm}$\RightarrowA$=$4\,\mathrm{cm}$.

Vì t = 0 lúc vật qua vị trí cân bằng theo chiều dương \Rightarrow  pha ban đầu là: $\varphi = - \dfrac{\pi}{2}$ (Hình I.2Gb)

Vậy phương trình dao động của vật là: x = 4cos $\left( {\pi t - \dfrac{\pi}{2}} \right)$ (cm).
}
\end{ex}

% ===== Câu 7 =====
% ID: L11C1B1-35-TLN
% Mức: VD
\dangbt{quãng đường}
\begin{ex}
% ID: L11C1B1-35-TLN
% Mức: VD
L11C1 Ôn Tập chương dao động – quãng đường: Vật có biên độ $10\,\text{cm}$. Quãng đường trong 3 chu kì theo cm bằng bao nhiêu?
\shortans{120}
\end{ex}

% ===== Câu 8 =====
% ID: L11C1B1-36-TN
% Mức: NB
\dangbt{Chưa có dạng}
\begin{ex}
% ID: L11C1B1-36-TN
% Mức: NB
Một con lắc lò xo gồm một lò xo có độ cứng không đổi. Khi khối lượng quả nặng là m thì tần số dao động là $1 \mathrm{Hz}$. Khi khối lượng quả nặng là $2\,\text{m}$ thì tần số dao động của con lắc là 1 $\nguon{ly11 kntt.pdf}$
\choice
{$2 \mathrm{Hz}$.}
{$\sqrt{2} \mathrm{Hz}.$}
{\True $\mathrm{Hz}.$}
{$0,5 \mathrm{Hz}$. $\sqrt{2}$}
\loigiai{
Chọn C.

Do f $. \sim\dfrac{1}{\sqrt{m}}\Rightarrow\dfrac{f_{1}}{f_{2}} = \sqrt{\dfrac{m_{2}}{m_{1}}}\Rightarrow\dfrac{1}{f_{2}} = \sqrt{\dfrac{$ 2 $\mathrm{m}$ }{m}}\Rightarrow f_{2} = \dfrac{1}{\sqrt{2}}Hz \right.
}
\end{ex}

% ===== Câu 9 =====
% ID: L11C1B1-37-DS
% Mức: VD
\dangbt{quãng đường}
\begin{ex}
% ID: L11C1B1-37-DS
% Mức: VD
L11C1 Ôn Tập chương dao động – quãng đường: Vật có biên độ $8\,\text{cm}$, chu kì $4\,\text{s}$.
\choiceTF
{\True Một chu kì đi $32\,\text{cm}$.}
{\True Nửa chu kì đi $16\,\text{cm}$.}
{\True Sau một chu kì độ dịch chuyển bằng 0.}
{\True Tốc độ trung bình là $8\,\text{cm}$ /s.}
\end{ex}

% ===== Câu 10 =====
% ID: L11C1B1-38-DS
% Mức: TH
\dangbt{Chưa có dạng}
\begin{ex}
% ID: L11C1B1-38-DS
% Mức: TH
Một con lắc lò xo gồm lò xo có độ cứng $k = 160\$, $\mathrm{N}/m$ và vật nặng có khối lượng $m = 400\$, $\mathrm{g},$ đặt trên mặt phẳng nằm ngang. Hệ số ma sát trượt giữa vật và mặt phẳng nằm ngang là $\mu = 0$,0005. Lấy $g =$ 10 $\mathrm{m/s}$ 2 $. Kéo vật lệch khỏi vị trí lò xo không biến dạng một đoạn$ 5 \mathrm{cm} $(theo phương của trục lò xo). Tại$ t = 0 $, buông nhẹ để vật dao động. Tính thời gian kể từ lúc vật bắt đầu dao động cho đến khi vật dừng hẳn.$ \nguon{ly11kntt.pdf}
\choiceTF

\loigiai{
Giả sử biên độ dao động của vật ở $\dfrac{1}{4}$ chu kì là A₁, biên độ dao động của vật ở $\dfrac{1}{4}$ chu kì tiếp theo là A₂. Công của lực ma sát gây ra độ biến đổi cơ năng của vật nặng trong nửa chu kì đầu là:

$\left. - \mu\text{mg}\left( {\text{A}_{1} + \text{A}_{2}} \right) = \dfrac{\text{kA}_{2}^{2}}{2} - \dfrac{k\text{A}_{1}^{2}}{2}\Rightarrow\text{A}_{2} - \text{A}_{1} = - \dfrac{2\mu\text{mg}}{\text{k}} = - \dfrac{2.0,0005.0,4.10}{160} \right.$

= -2,5$\cdot$ 10^{-5} m = -$0,0025\,\mathrm{cm}$

Số nửa chu kì vật thực hiện được từ t = 0 đến khi dừng hẳn là:

$\text{N} = \dfrac{\text{A}_{1}}{0,0025} = \dfrac{5}{0,0025} = 2000$ lần

Cứ mỗi chu kì có hai lần vật qua vị trí đó. Vậy thời gian để vật qua vị trí đó 2000 lần bằng 1000 lần chu kì: $\text{\Delta t} = 1000\text{T} = 1000.2\pi\sqrt{\dfrac{\text{m}}{\text{k}}} = 1000.2\pi\sqrt{\dfrac{0,4}{160}} = 314\text{s}$.
}
\end{ex}

% ===== Câu 11 =====
% ID: L11C1B1-39-TLN
% Mức: VD
\dangbt{quãng đường}
\begin{ex}
% ID: L11C1B1-39-TLN
% Mức: VD
L11C1 Ôn Tập chương dao động – quãng đường: Vật có biên độ $5\,\text{cm}$. Quãng đường trong 3 chu kì theo cm bằng bao nhiêu?
\shortans{60}
\end{ex}

% ===== Câu 12 =====
% ID: L11C1B1-40-TN
% Mức: NB
\dangbt{Chưa có dạng}
\begin{ex}
% ID: L11C1B1-40-TN
% Mức: NB
Lợi ích của hiện tượng cộng hưởng được ứng dụng trong trường hợp nào sau đây? $\nguon{ly11 kntt.pdf}$
\choice
{\True Chế tạo máy phát tần số.}
{Chế tạo bộ phận giảm xóc của ô tô, xe máy.}
{Lắp đặt các động cơ điện trong nhà xưởng.}
{Thiết kế các công trình ở những vùng thường có địa chấn.}
\loigiai{
Chọn A.
}
\end{ex}

% ===== Câu 13 =====
% ID: L11C1B1-41-DS
% Mức: TH
\dangbt{Thời điểm vật qua x1}
\begin{ex}
% ID: L11C1B1-41-DS
% Mức: TH
L11C1 Ôn Tập chương dao động – Thời điểm vật qua x1: Vật có biên độ $9\,\text{cm}$, chu kì $1\,\text{s}$.
\choiceTF
{\True Một chu kì đi $36\,\text{cm}$.}
{\True Nửa chu kì đi $18\,\text{cm}$.}
{\True Sau một chu kì độ dịch chuyển bằng 0.}
{\True Tốc độ trung bình là $36\,\text{cm}$ /s.}
\end{ex}

% ===== Câu 14 =====
% ID: L11C1B1-42-DS
% Mức: TH
\dangbt{quãng đường}
\begin{ex}
% ID: L11C1B1-42-DS
% Mức: TH
L11C1 Ôn Tập chương dao động – quãng đường: Vật có biên độ $6\,\text{cm}$, chu kì $1\,\text{s}$. Tính quãng đường và tốc độ trung bình trong 3 chu kì.
\choiceTF

\end{ex}

% ===== Câu 15 =====
% ID: L11C1B1-43-TLN
% Mức: TH
\dangbt{Thời điểm vật qua x1}
\begin{ex}
% ID: L11C1B1-43-TLN
% Mức: TH
L11C1 Ôn Tập chương dao động – Thời điểm vật qua x1: Vật có biên độ $6\,\text{cm}$. Quãng đường trong 3 chu kì theo cm bằng bao nhiêu?
\shortans{72}
\end{ex}

% ===== Câu 16 =====
% ID: L11C1B1-44-TN
% Mức: TH
\dangbt{Chưa có dạng}
\begin{ex}
% ID: L11C1B1-44-TN
% Mức: TH
Tìm phát biểu sai về gia tốc của một vật dao động điều hoà. $\nguon{ly11 kntt.pdf}$
\choice
{Gia tốc đổi chiều khi vật đi qua vị trí cân bằng.}
{\True Gia tốc luôn ngược chiều với vận tốc.}
{Gia tốc luôn hướng về vị trí cân bằng.}
{Gia tốc biến đổi ngược pha với li độ.}
\loigiai{
Chọn B.

Có những giai đoạn gia tốc cùng chiều với vận tốc, đó là khi vật đi từ vị trí biên về vị trí cân bằng.
}
\end{ex}

% ===== Câu 17 =====
% ID: L11C1B1-45-DS
% Mức: VD
\dangbt{Thời điểm vật qua x1}
\begin{ex}
% ID: L11C1B1-45-DS
% Mức: VD
L11C1 Ôn Tập chương dao động – Thời điểm vật qua x1: Vật có biên độ $10\,\text{cm}$, chu kì $2\,\text{s}$.
\choiceTF
{\True Một chu kì đi $40\,\text{cm}$.}
{\True Nửa chu kì đi $20\,\text{cm}$.}
{\True Sau một chu kì độ dịch chuyển bằng 0.}
{\True Tốc độ trung bình là $20\,\text{cm}$ /s.}
\end{ex}

% ===== Câu 18 =====
% ID: L11C1B1-46-DS
% Mức: VD
\dangbt{quãng đường}
\begin{ex}
% ID: L11C1B1-46-DS
% Mức: VD
L11C1 Ôn Tập chương dao động – quãng đường: Vật có biên độ $7\,\text{cm}$, chu kì $2\,\text{s}$. Tính quãng đường và tốc độ trung bình trong 3 chu kì.
\choiceTF

\end{ex}

% ===== Câu 19 =====
% ID: L11C1B1-47-TLN
% Mức: VD
\dangbt{Thời điểm vật qua x1}
\begin{ex}
% ID: L11C1B1-47-TLN
% Mức: VD
L11C1 Ôn Tập chương dao động – Thời điểm vật qua x1: Vật có biên độ $7\,\text{cm}$. Quãng đường trong 3 chu kì theo cm bằng bao nhiêu?
\shortans{84}
\end{ex}

% ===== Câu 20 =====
% ID: L11C1B1-48-TN
% Mức: TH
\dangbt{Chưa có dạng}
\begin{ex}
% ID: L11C1B1-48-TN
% Mức: TH
Một con lắc lò xo nằm ngang, đang thực hiện dao động điều hoà. Tìm phát biểu sai. $\nguon{ly11 kntt.pdf}$
\choice
{Động năng của vật nặng và thế năng đàn hồi của lò xo là hai thành phần tạo thành cơ năng của con lắc.}
{Động năng và thế năng của con lắc biến thiên tuần hoàn với cùng một tần số như nhau.}
{Khi vật ở một trong hai vị trí biên thì thế năng của con lắc đạt giá trị cực đại.}
{\True Động năng và thế năng của con lắc biến thiên tuần hoàn với cùng chu kì như chu kì của dao động.}
\loigiai{
Chọn D.

Chu kì biến thiên của động năng và thế năng bằng nửa chu kì của dao động.
}
\end{ex}

% ===== Câu 21 =====
% ID: L11C1B1-49-DS
% Mức: VD
\dangbt{Thời điểm vật qua x1}
\begin{ex}
% ID: L11C1B1-49-DS
% Mức: VD
L11C1 Ôn Tập chương dao động – Thời điểm vật qua x1: Vật có biên độ $5\,\text{cm}$, chu kì $4\,\text{s}$.
\choiceTF
{\True Một chu kì đi $20\,\text{cm}$.}
{\True Nửa chu kì đi $10\,\text{cm}$.}
{\True Sau một chu kì độ dịch chuyển bằng 0.}
{\True Tốc độ trung bình là $5\,\text{cm}$ /s.}
\end{ex}

% ===== Câu 22 =====
% ID: L11C1B1-50-DS
% Mức: VD
\dangbt{quãng đường}
\begin{ex}
% ID: L11C1B1-50-DS
% Mức: VD
L11C1 Ôn Tập chương dao động – quãng đường: Vật có biên độ $8\,\text{cm}$, chu kì $4\,\text{s}$. Tính quãng đường và tốc độ trung bình trong 3 chu kì.
\choiceTF

\end{ex}

% ===== Câu 23 =====
% ID: L11C1B1-51-TLN
% Mức: VD
\dangbt{Thời điểm vật qua x1}
\begin{ex}
% ID: L11C1B1-51-TLN
% Mức: VD
L11C1 Ôn Tập chương dao động – Thời điểm vật qua x1: Vật có biên độ $8\,\text{cm}$. Quãng đường trong 3 chu kì theo cm bằng bao nhiêu?
\shortans{96}
\end{ex}

% ===== Câu 24 =====
% ID: L11C1B1-52-TN
% Mức: TH
\dangbt{Chưa có dạng}
\begin{ex}
% ID: L11C1B1-52-TN
% Mức: TH
Tìm phát biễu sai về dao động tắt dần của con lắc lò xo. $\nguon{ly11 kntt.pdf}$
\choice
{Cơ năng của con lắc luôn giảm dần.}
{Động năng của vật có lúc tăng, lúc giảm.}
{\True Động năng của vật luôn giảm dần.}
{Thế năng của con lắc có lúc tăng, lúc giảm.}
\loigiai{
Chọn C.

Trong dao động tắt dần, vẫn có những lúc động năng tăng, đó là lúc vật đi từ vị trí biên về vị trí cân bằng.
}
\end{ex}

% ===== Câu 25 =====
% ID: L11C1B1-53-DS
% Mức: TH
\dangbt{Viết phương trình dao động điều hòa.}
\begin{ex}
% ID: L11C1B1-53-DS
% Mức: TH
L11C1 Ôn Tập chương dao động – Viết phương trình dao động điều hòa.: Vật có $x=6\cos(2\pi t/1+\pi/3)$ cm.
\choiceTF
{\True Biên độ $6\,\text{cm}$.}
{\True Chu kì $1\,\text{s}$.}
{\True Li độ đầu $3\,\text{cm}$.}
{Ban đầu vật đi theo chiều dương.}
\end{ex}

% ===== Câu 26 =====
% ID: L11C1B1-54-DS
% Mức: TH
\dangbt{Thời điểm vật qua x1}
\begin{ex}
% ID: L11C1B1-54-DS
% Mức: TH
L11C1 Ôn Tập chương dao động – Thời điểm vật qua x1: Vật có biên độ $9\,\text{cm}$, chu kì $1\,\text{s}$. Tính quãng đường và tốc độ trung bình trong 3 chu kì.
\choiceTF

\end{ex}

% ===== Câu 27 =====
% ID: L11C1B1-55-TLN
% Mức: TH
\dangbt{Viết phương trình dao động điều hòa.}
\begin{ex}
% ID: L11C1B1-55-TLN
% Mức: TH
L11C1 Ôn Tập chương dao động – Viết phương trình dao động điều hòa.: Vật có biên độ $9\,\text{cm}$. Chiều dài quỹ đạo theo cm bằng bao nhiêu?
\shortans{18}
\end{ex}

% ===== Câu 28 =====
% ID: L11C1B1-56-TN
% Mức: VD
\dangbt{quãng đường}
\begin{ex}
% ID: L11C1B1-56-TN
% Mức: VD
Một vật dao động điều hòa với biên độ $A$ và chu kì $T$. Thời gian ngắn nhất để vật đi được quãng đường có độ dài $A\sqrt{2}$ là
\choice
{$\dfrac{T}{8}$}
{\True $\dfrac{T}{4}$}
{$\dfrac{T}{6}$}
{$\dfrac{T}{12}$}
\loigiai{
Thời gian ngắn nhất ứng với quãng đường lớn nhất trong khoảng thời gian đó:
$S=2A\sin\left(\dfrac{\pi\Delta t}{T}\right).$
Theo đề:
$2A\sin\left(\dfrac{\pi\Delta t}{T}\right)=A\sqrt{2}.$
Suy ra
$\sin\left(\dfrac{\pi\Delta t}{T}\right)=\dfrac{\sqrt{2}}{2}.$
Do đó
 $\dfrac{\pi\Delta t}{T}=\dfrac{\pi}{4}\Rightarrow\Delta$ t= $\dfrac{T}{4}$.
}
\end{ex}

% ===== Câu 29 =====
% ID: L11C1B1-57-DS
% Mức: VD
\dangbt{Viết phương trình dao động điều hòa.}
\begin{ex}
% ID: L11C1B1-57-DS
% Mức: VD
L11C1 Ôn Tập chương dao động – Viết phương trình dao động điều hòa.: Vật có $x=7\cos(2\pi t/2+\pi/3)$ cm.
\choiceTF
{\True Biên độ $7\,\text{cm}$.}
{\True Chu kì $2\,\text{s}$.}
{\True Li độ đầu $3.5\,\text{cm}$.}
{Ban đầu vật đi theo chiều dương.}
\end{ex}

% ===== Câu 30 =====
% ID: L11C1B1-58-DS
% Mức: VD
\dangbt{Thời điểm vật qua x1}
\begin{ex}
% ID: L11C1B1-58-DS
% Mức: VD
L11C1 Ôn Tập chương dao động – Thời điểm vật qua x1: Vật có biên độ $10\,\text{cm}$, chu kì $2\,\text{s}$. Tính quãng đường và tốc độ trung bình trong 3 chu kì.
\choiceTF

\end{ex}

% ===== Câu 31 =====
% ID: L11C1B1-59-TLN
% Mức: VD
\dangbt{Viết phương trình dao động điều hòa.}
\begin{ex}
% ID: L11C1B1-59-TLN
% Mức: VD
L11C1 Ôn Tập chương dao động – Viết phương trình dao động điều hòa.: Vật có biên độ $10\,\text{cm}$. Chiều dài quỹ đạo theo cm bằng bao nhiêu?
\shortans{20}
\end{ex}

% ===== Câu 32 =====
% ID: L11C1B1-60-TN
% Mức: VD
\dangbt{quãng đường}
\begin{ex}
% ID: L11C1B1-60-TN
% Mức: VD
Một vật dao động điều hòa với phương trình
$x=4\cos\left(4\pi t+\dfrac{\pi}{3}\right)\ \text{cm},$
với $t$ đo bằng giây. Quãng đường lớn nhất mà vật đi được trong khoảng thời gian $\dfrac{1}{6}\ \text{s}$ là
\choice
{$\sqrt{3}\ \text{cm}$}
{$3\sqrt{3}\ \text{cm}$}
{$2\sqrt{3}\ \text{cm}$}
{\True $4\sqrt{3}\ \text{cm}$}
\loigiai{
Từ phương trình, ta có
$A=4\ \mathrm{cm},\qquad \omega=4\pi\ \mathrm{rad/s}.$
Suy ra
$T=\dfrac{2\pi}{\omega}=\dfrac{2\pi}{4\pi}=\dfrac{1}{2}\ \mathrm{s}.$
Với $\Delta t=\dfrac{1}{6}\ \mathrm{s}$:
$\dfrac{\Delta t}{T}=\dfrac{1/6}{1/2}=\dfrac{1}{3}.$
Do đó
$S_{\max}=2A\sin\left(\dfrac{\pi\Delta t}{T}\right)=2\cdot 4\sin\dfrac{\pi}{3}=4\sqrt{3}\ \mathrm{cm}.$
}
\end{ex}

% ===== Câu 33 =====
% ID: L11C1B1-61-DS
% Mức: VD
\dangbt{Viết phương trình dao động điều hòa.}
\begin{ex}
% ID: L11C1B1-61-DS
% Mức: VD
L11C1 Ôn Tập chương dao động – Viết phương trình dao động điều hòa.: Vật có $x=8\cos(2\pi t/4+\pi/3)$ cm.
\choiceTF
{\True Biên độ $8\,\text{cm}$.}
{\True Chu kì $4\,\text{s}$.}
{\True Li độ đầu $4\,\text{cm}$.}
{Ban đầu vật đi theo chiều dương.}
\end{ex}

% ===== Câu 34 =====
% ID: L11C1B1-62-DS
% Mức: VD
\dangbt{Thời điểm vật qua x1}
\begin{ex}
% ID: L11C1B1-62-DS
% Mức: VD
L11C1 Ôn Tập chương dao động – Thời điểm vật qua x1: Vật có biên độ $5\,\text{cm}$, chu kì $4\,\text{s}$. Tính quãng đường và tốc độ trung bình trong 3 chu kì.
\choiceTF

\end{ex}

% ===== Câu 35 =====
% ID: L11C1B1-63-TLN
% Mức: VD
\dangbt{Viết phương trình dao động điều hòa.}
\begin{ex}
% ID: L11C1B1-63-TLN
% Mức: VD
L11C1 Ôn Tập chương dao động – Viết phương trình dao động điều hòa.: Vật có biên độ $5\,\text{cm}$. Chiều dài quỹ đạo theo cm bằng bao nhiêu?
\shortans{10}
\end{ex}

% ===== Câu 36 =====
% ID: L11C1B1-64-TN
% Mức: VD
\dangbt{quãng đường}
\begin{ex}
% ID: L11C1B1-64-TN
% Mức: VD
Một vật dao động điều hòa với chu kì $T$ và biên độ $A$. Quãng đường vật đi được tối đa trong khoảng thời gian $\dfrac{7T}{6}$ là
\choice
{\True $5A$}
{$A$}
{$A\sqrt{3}$}
{$1,5A\sqrt{3}$}
\loigiai{
Ta có
$\dfrac{7T}{6}=T+\dfrac{T}{6}.$
Trong một chu kì, vật đi được quãng đường $4A$.
Trong thời gian $\dfrac{T}{6}$, quãng đường lớn nhất là
$2A\sin\dfrac{\pi}{6}=A.$
Vậy
$S_{\max}=4A+A=5A.$
}
\end{ex}

% ===== Câu 37 =====
% ID: L11C1B1-65-DS
% Mức: TH
\dangbt{vận tốc trung bình và tốc độ trung bình}
\begin{ex}
% ID: L11C1B1-65-DS
% Mức: TH
Ôn tập – vận tốc trung bình và tốc độ trung bình: Vật đi từ $x=2$ cm đến $x=-4$ cm trong $0,5\,\text{s}$, quãng đường đi được $10\,\text{cm}$.
\choiceTF
{\True Độ dịch chuyển bằng $-6$ cm.}
{\True Vận tốc trung bình bằng $-12$ cm/s.}
{\True Tốc độ trung bình bằng $20\,\text{cm}$ /s.}
{Độ lớn vận tốc trung bình bằng tốc độ trung bình.}
\end{ex}

% ===== Câu 38 =====
% ID: L11C1B1-66-TL
% Mức: TH
\dangbt{Viết phương trình dao động điều hòa.}
\begin{bt}
% ID: L11C1B1-66-TL
% Mức: TH
L11C1 Ôn Tập chương dao động – Viết phương trình dao động điều hòa.: Vật có biên độ $6\,\text{cm}$, chu kì $1\,\text{s}$, ban đầu ở biên dương. Viết phương trình dao động.
\end{bt}

% ===== Câu 39 =====
% ID: L11C1B1-67-TLN
% Mức: TH
\dangbt{vận tốc trung bình và tốc độ trung bình}
\begin{ex}
% ID: L11C1B1-67-TLN
% Mức: TH
Ôn tập – vận tốc trung bình và tốc độ trung bình: Vật đi quãng đường $24\,\text{cm}$ trong $1,5\,\text{s}$. Tốc độ trung bình theo cm/s bằng bao nhiêu?
\shortans{16}
\end{ex}

% ===== Câu 40 =====
% ID: L11C1B1-68-TN
% Mức: VD
\dangbt{quãng đường}
\begin{ex}
% ID: L11C1B1-68-TN
% Mức: VD
Một vật dao động điều hòa với biên độ $A$ và chu kì $T$. Thời gian dài nhất để vật đi được quãng đường có độ dài $A$ là
\choice
{$\dfrac{T}{6}$}
{$\dfrac{T}{4}$}
{\True $\dfrac{T}{3}$}
{$\dfrac{T}{8}$}
\loigiai{
Thời gian dài nhất để đi được quãng đường $A$ xảy ra khi vật chuyển động gần biên.
Khi đó
$S=2A\left[1-\cos\left(\dfrac{\pi\Delta t}{T}\right)\right].$
Cho $S=A$:
$2A\left[1-\cos\left(\dfrac{\pi\Delta t}{T}\right)\right]=A.$
Suy ra
$\cos\left(\dfrac{\pi\Delta t}{T}\right)=\dfrac{1}{2}.$
Do đó
 $\dfrac{\pi\Delta t}{T}=\dfrac{\pi}{3}\Rightarrow\Delta$ t= $\dfrac{T}{3}$.
}
\end{ex}

% ===== Câu 41 =====
% ID: L11C1B1-69-DS
% Mức: VD
\dangbt{vận tốc trung bình và tốc độ trung bình}
\begin{ex}
% ID: L11C1B1-69-DS
% Mức: VD
Ôn tập – vận tốc trung bình và tốc độ trung bình: Trong một chu kì, vật dao động điều hòa biên độ $5\,\text{cm}$, chu kì $2\,\text{s}$.
\choiceTF
{\True Độ dịch chuyển bằng 0.}
{\True Vận tốc trung bình bằng 0.}
{\True Quãng đường bằng $20\,\text{cm}$.}
{Tốc độ trung bình bằng $5\,\text{cm}$ /s.}
\end{ex}

% ===== Câu 42 =====
% ID: L11C1B1-70-TL
% Mức: VD
\dangbt{Viết phương trình dao động điều hòa.}
\begin{bt}
% ID: L11C1B1-70-TL
% Mức: VD
L11C1 Ôn Tập chương dao động – Viết phương trình dao động điều hòa.: Vật có biên độ $7\,\text{cm}$, chu kì $2\,\text{s}$, ban đầu ở biên dương. Viết phương trình dao động.
\end{bt}

% ===== Câu 43 =====
% ID: L11C1B1-71-TLN
% Mức: VD
\dangbt{vận tốc trung bình và tốc độ trung bình}
\begin{ex}
% ID: L11C1B1-71-TLN
% Mức: VD
Ôn tập – vận tốc trung bình và tốc độ trung bình: Vật chuyển từ $x=-3$ cm đến $x=5$ cm trong $0,4\,\text{s}$. Vận tốc trung bình theo cm/s bằng bao nhiêu?
\shortans{20}
\end{ex}

% ===== Câu 44 =====
% ID: L11C1B1-72-TN
% Mức: VD
\dangbt{quãng đường}
\begin{ex}
% ID: L11C1B1-72-TN
% Mức: VD
Một vật dao động điều hòa trên trục $Ox$. Gọi $t_1$ và $t_2$ lần lượt là khoảng thời gian ngắn nhất và dài nhất để vật đi được quãng đường bằng biên độ. Tỉ số $\dfrac{t_1}{t_2}$ bằng
\choice
{$2$}
{\True $\dfrac{1}{2}$}
{$\dfrac{1}{3}$}
{$0,5\sqrt{2}$}
\loigiai{
Thời gian ngắn nhất để vật đi được quãng đường $A$ là
$t_1=\dfrac{T}{6}.$
Thời gian dài nhất để vật đi được quãng đường $A$ là
$t_2=\dfrac{T}{3}.$
Do đó
 $\dfrac{t_1}{t_2}$ = $\dfrac{T/6}{T/3}$ = $\dfrac{1}{2}$.
}
\end{ex}

% ===== Câu 45 =====
% ID: L11C1B1-73-DS
% Mức: VD
\dangbt{vận tốc trung bình và tốc độ trung bình}
\begin{ex}
% ID: L11C1B1-73-DS
% Mức: VD
Ôn tập – vận tốc trung bình và tốc độ trung bình: Vật đi từ biên dương đến biên âm trong nửa chu kì.
\choiceTF
{\True Độ dịch chuyển bằng $-2A$.}
{\True Quãng đường bằng $2A$.}
{\True Độ lớn vận tốc trung bình bằng tốc độ trung bình.}
{Vận tốc trung bình luôn dương.}
\end{ex}

% ===== Câu 46 =====
% ID: L11C1B1-74-TL
% Mức: VD
\dangbt{Viết phương trình dao động điều hòa.}
\begin{bt}
% ID: L11C1B1-74-TL
% Mức: VD
L11C1 Ôn Tập chương dao động – Viết phương trình dao động điều hòa.: Vật có biên độ $8\,\text{cm}$, chu kì $4\,\text{s}$, ban đầu ở biên dương. Viết phương trình dao động.
\end{bt}

% ===== Câu 47 =====
% ID: L11C1B1-75-TLN
% Mức: VD
\dangbt{vận tốc trung bình và tốc độ trung bình}
\begin{ex}
% ID: L11C1B1-75-TLN
% Mức: VD
Ôn tập – vận tốc trung bình và tốc độ trung bình: Vật dao động biên độ $6\,\text{cm}$, chu kì $3\,\text{s}$. Tốc độ trung bình trong một chu kì theo cm/s bằng bao nhiêu?
\shortans{8}
\end{ex}

% ===== Câu 48 =====
% ID: L11C1B1-76-TN
% Mức: VDC
\dangbt{quãng đường}
\begin{ex}
% ID: L11C1B1-76-TN
% Mức: VDC
Một chất điểm dao động điều hòa theo phương trình
$x=5\cos 4\pi t\ \text{cm},$
với $t$ đo bằng giây. Trong khoảng thời gian $\dfrac{7}{6}\ \text{s}$, quãng đường lớn nhất vật có thể đi được là
\choice
{$42,5\ \text{cm}$}
{\True $48,66\ \text{cm}$}
{$45\ \text{cm}$}
{$30\sqrt{3}\ \text{cm}$}
\loigiai{
Từ phương trình:
$A=5\ \text{cm},\qquad \omega=4\pi\ \text{rad/s}.$
Suy ra
$T=\dfrac{2\pi}{4\pi}=0,5\ \text{s}.$
Ta có
$\dfrac{7}{6}\ \text{s}=2T+\dfrac{T}{3}.$
Trong $2T$, vật đi được
$2\cdot 4A=8A=40\ \text{cm}.$
Trong thời gian $\dfrac{T}{3}$, quãng đường lớn nhất là
$A\sqrt{3}=5\sqrt{3}\approx 8,66\ \text{cm}.$
Vậy
$S_{\max}=40+8,66=48,66\ \text{cm}.$
}
\end{ex}

% ===== Câu 49 =====
% ID: L11C1B1-77-DS
% Mức: TH
\dangbt{Xác định và tính toán các đại lượng vận tốc, gia tốc tại thời điểm xác định}
\begin{ex}
% ID: L11C1B1-77-DS
% Mức: TH
L11C1 Ôn Tập chương dao động – Xác định và tính toán các đại lượng vận tốc, gia tốc tại thời điểm xác định: Vật dao động theo $x=7\cos(4t)$ cm.
\choiceTF
{\True $v=-28\sin(4t)$ cm/s.}
{\True $a=-112\cos(4t)$ cm/s².}
{\True Tốc độ cực đại $28\,\text{cm}$ /s.}
{Gia tốc bằng 0 ở biên.}
\end{ex}

% ===== Câu 50 =====
% ID: L11C1B1-78-DS
% Mức: TH
\dangbt{vận tốc trung bình và tốc độ trung bình}
\begin{ex}
% ID: L11C1B1-78-DS
% Mức: TH
Ôn tập – vận tốc trung bình và tốc độ trung bình: Vật đi từ $x=4$ cm đến $x=-2$ cm trong $0,5\,\text{s}$ và đi được $10\,\text{cm}$. Tính vận tốc trung bình và tốc độ trung bình.
\choiceTF

\end{ex}

% ===== Câu 51 =====
% ID: L11C1B1-79-TLN
% Mức: TH
\dangbt{Xác định và tính toán các đại lượng vận tốc, gia tốc tại thời điểm xác định}
\begin{ex}
% ID: L11C1B1-79-TLN
% Mức: TH
L11C1 Ôn Tập chương dao động – Xác định và tính toán các đại lượng vận tốc, gia tốc tại thời điểm xác định: Vật có $A=10$ cm, $\omega=3$ rad/s. Tốc độ cực đại theo cm/s bằng bao nhiêu?
\shortans{30}
\end{ex}

% ===== Câu 52 =====
% ID: L11C1B1-80-TN
% Mức: VDC
\dangbt{quãng đường}
\begin{ex}
% ID: L11C1B1-80-TN
% Mức: VDC
Một vật dao động điều hòa với biên độ $10\ \text{cm}$. Quãng đường nhỏ nhất mà vật đi được trong $0,5\ \text{s}$ là $10\ \text{cm}$. Tốc độ lớn nhất của vật là
\choice
{$39,95\ \text{cm/s}$}
{\True $41,9\ \text{cm/s}$}
{$40,65\ \text{cm/s}$}
{$41,2\ \text{cm/s}$}
\loigiai{
Ta có
$A=10\ \mathrm{cm},\qquad \Delta t=0,5\ \mathrm{s}.$
Theo đề:
$S_{\min}=10\ \mathrm{cm}=A.$
Mà
$S_{\min}=2A\left[1-\cos\left(\dfrac{\pi\Delta t}{T}\right)\right].$
Suy ra
$2A\left[1-\cos\left(\dfrac{\pi\Delta t}{T}\right)\right]=A \Rightarrow \cos\left(\dfrac{\pi\Delta t}{T}\right)=\dfrac{1}{2}.$
Do đó $\dfrac{\pi\Delta t}{T}=\dfrac{\pi}{3} \Rightarrow T=3\Delta t=1,5\ \mathrm{s}.$
Tốc độ cực đại:
$v_{\max}=\omega A=\dfrac{2\pi}{T}A=\dfrac{2\pi}{1,5}\cdot 10 \approx 41,9\ \mathrm{cm/s}.$
}
\end{ex}

% ===== Câu 53 =====
% ID: L11C1B1-81-DS
% Mức: VD
\dangbt{Xác định và tính toán các đại lượng vận tốc, gia tốc tại thời điểm xác định}
\begin{ex}
% ID: L11C1B1-81-DS
% Mức: VD
L11C1 Ôn Tập chương dao động – Xác định và tính toán các đại lượng vận tốc, gia tốc tại thời điểm xác định: Vật dao động theo $x=8\cos(5t)$ cm.
\choiceTF
{\True $v=-40\sin(5t)$ cm/s.}
{\True $a=-200\cos(5t)$ cm/s².}
{\True Tốc độ cực đại $40\,\text{cm}$ /s.}
{Gia tốc bằng 0 ở biên.}
\end{ex}

% ===== Câu 54 =====
% ID: L11C1B1-82-DS
% Mức: VD
\dangbt{vận tốc trung bình và tốc độ trung bình}
\begin{ex}
% ID: L11C1B1-82-DS
% Mức: VD
Ôn tập – vận tốc trung bình và tốc độ trung bình: Vật dao động điều hòa biên độ $8\,\text{cm}$, chu kì $2\,\text{s}$. Tính vận tốc trung bình và tốc độ trung bình trong một chu kì.
\choiceTF

\end{ex}

% ===== Câu 55 =====
% ID: L11C1B1-83-TLN
% Mức: VD
\dangbt{Xác định và tính toán các đại lượng vận tốc, gia tốc tại thời điểm xác định}
\begin{ex}
% ID: L11C1B1-83-TLN
% Mức: VD
L11C1 Ôn Tập chương dao động – Xác định và tính toán các đại lượng vận tốc, gia tốc tại thời điểm xác định: Vật có $A=5$ cm, $\omega=4$ rad/s. Tốc độ cực đại theo cm/s bằng bao nhiêu?
\shortans{20}
\end{ex}

% ===== Câu 56 =====
% ID: L11C1B1-84-TN
% Mức: VDC
\dangbt{quãng đường}
\begin{ex}
% ID: L11C1B1-84-TN
% Mức: VDC
Một vật dao động điều hòa dọc theo trục $Ox$, quanh vị trí cân bằng $O$ với biên độ $A$ và chu kì $T$. Trong khoảng thời gian $\dfrac{T}{3}$, quãng đường lớn nhất mà vật có thể đi được là
\choice
{$A$}
{$1{,}5A$}
{\True $A\sqrt{3}$}
{$A\sqrt{2}$}
\loigiai{
Với $0\lt \Delta t<\dfrac{T}{2}$:
$S_{\max}=2A\sin\left(\dfrac{\pi\Delta t}{T}\right).$
Thay $\Delta t=\dfrac{T}{3}$:
$S_{\max}=2A\sin\dfrac{\pi}{3}=A\sqrt{3}.$
}
\end{ex}

% ===== Câu 57 =====
% ID: L11C1B1-85-DS
% Mức: VD
\dangbt{Xác định và tính toán các đại lượng vận tốc, gia tốc tại thời điểm xác định}
\begin{ex}
% ID: L11C1B1-85-DS
% Mức: VD
L11C1 Ôn Tập chương dao động – Xác định và tính toán các đại lượng vận tốc, gia tốc tại thời điểm xác định: Vật dao động theo $x=9\cos(2t)$ cm.
\choiceTF
{\True $v=-18\sin(2t)$ cm/s.}
{\True $a=-36\cos(2t)$ cm/s².}
{\True Tốc độ cực đại $18\,\text{cm}$ /s.}
{Gia tốc bằng 0 ở biên.}
\end{ex}

% ===== Câu 58 =====
% ID: L11C1B1-86-DS
% Mức: VD
\dangbt{vận tốc trung bình và tốc độ trung bình}
\begin{ex}
% ID: L11C1B1-86-DS
% Mức: VD
Ôn tập – vận tốc trung bình và tốc độ trung bình: Vật đi từ biên dương đến biên âm trong thời gian $T/2$. So sánh độ lớn vận tốc trung bình với tốc độ trung bình.
\choiceTF

\end{ex}

% ===== Câu 59 =====
% ID: L11C1B1-87-TLN
% Mức: VD
\dangbt{Xác định và tính toán các đại lượng vận tốc, gia tốc tại thời điểm xác định}
\begin{ex}
% ID: L11C1B1-87-TLN
% Mức: VD
L11C1 Ôn Tập chương dao động – Xác định và tính toán các đại lượng vận tốc, gia tốc tại thời điểm xác định: Vật có $A=6$ cm, $\omega=5$ rad/s. Tốc độ cực đại theo cm/s bằng bao nhiêu?
\shortans{30}
\end{ex}

% ===== Câu 60 =====
% ID: L11C1B1-88-TN
% Mức: VDC
\dangbt{quãng đường}
\begin{ex}
% ID: L11C1B1-88-TN
% Mức: VDC
Một vật dao động điều hòa dọc theo trục $Ox$, quanh vị trí cân bằng $O$ với biên độ $A$ và chu kì $T$. Trong khoảng thời gian $\dfrac{T}{4}$, quãng đường nhỏ nhất mà vật có thể đi được là
\choice
{$(\sqrt{3}-1)A$}
{$1{,}5A$}
{$A\sqrt{3}$}
{\True $A(2-\sqrt{2})$}
\loigiai{
Với $0\lt \Delta t<\dfrac{T}{2}$, quãng đường nhỏ nhất vật đi được là
$S_{\min}=2A\left[1-\cos\left(\dfrac{\pi\Delta t}{T}\right)\right].$
Với $\Delta t=\dfrac{T}{4}$:
 S_{ $\min$ }=2 $A$ \left(1-\cos\dfrac{\pi}{4}\right) $=2$ A $\left(1-\dfrac{\sqrt{2}}{2}\right)$ = $A$ (2-\sqrt{2}).
}
\end{ex}

% ===== Câu 61 =====
% ID: L11C1B1-89-DS
% Mức: TH
\dangbt{Xác định và tính toán các đại lượng vận tốc, gia tốc tại thời điểm xác định}
\begin{ex}
% ID: L11C1B1-89-DS
% Mức: TH
L11C1 Ôn Tập chương dao động – Xác định và tính toán các đại lượng vận tốc, gia tốc tại thời điểm xác định: Từ $x=7\cos(2t+\pi/3)$ cm, hãy lập biểu thức vận tốc và gia tốc.
\choiceTF

\end{ex}

% ===== Câu 62 =====
% ID: L11C1B1-90-TN
% Mức: VDC
\dangbt{quãng đường}
\begin{ex}
% ID: L11C1B1-90-TN
% Mức: VDC
Chọn phương án sai khi nói về vật dao động điều hòa dọc theo trục $Ox$, với $O$ là vị trí cân bằng, biên độ $A$ và chu kì $T$.
\choice
{Thời gian ngắn nhất vật đi từ vị trí biên đến vị trí mà tại đó động năng bằng một nửa giá trị cực đại là $\dfrac{T}{8}$}
{Để đi được quãng đường $A$ cần thời gian tối thiểu là $\dfrac{T}{6}$}
{Quãng đường đi được tối thiểu trong khoảng thời gian $\dfrac{T}{3}$ là $A$}
{\True Thời gian ngắn nhất vật đi từ vị trí có li độ cực đại đến vị trí mà tại đó vật đi theo chiều dương đồng thời lực kéo về có độ lớn bằng nửa giá trị cực đại là $\dfrac{T}{6}$}
\loigiai{
Ta kiểm tra từng phát biểu.

Động năng bằng một nửa giá trị cực đại:
$W_{\text{đ}}=\dfrac{1}{2}W_{\text{đ max}}\Rightarrow|x|=\dfrac{A}{\sqrt{2}}$.
Từ biên đến $|x|=\dfrac{A}{\sqrt{2}}$ mất thời gian $\dfrac{T}{8}$, nên $A$ đúng.

Thời gian tối thiểu để đi được quãng đường $A$ là $\dfrac{T}{6}$, nên $B$ đúng.

Quãng đường nhỏ nhất trong thời gian $\dfrac{T}{3}$:
$S_{\min}=2A\left(1-\cos\dfrac{\pi}{3}\right)=A,$
nên $C$ đúng.

Phát biểu $D$ sai vì từ vị trí $x=+A$ đến vị trí mà vật đi theo chiều dương và lực kéo về có độ lớn bằng nửa cực đại không phải mất $\dfrac{T}{6}$.
}
\end{ex}

% ===== Câu 63 =====
% ID: L11C1B1-91-DS
% Mức: VD
\dangbt{Xác định và tính toán các đại lượng vận tốc, gia tốc tại thời điểm xác định}
\begin{ex}
% ID: L11C1B1-91-DS
% Mức: VD
L11C1 Ôn Tập chương dao động – Xác định và tính toán các đại lượng vận tốc, gia tốc tại thời điểm xác định: Từ $x=8\cos(3t+\pi/3)$ cm, hãy lập biểu thức vận tốc và gia tốc.
\choiceTF

\end{ex}

% ===== Câu 64 =====
% ID: L11C1B1-92-TN
% Mức: VDC
\dangbt{quãng đường}
\begin{ex}
% ID: L11C1B1-92-TN
% Mức: VDC
Cho vật dao động điều hòa biên độ $A$, chu kì $T$. Quãng đường lớn nhất mà vật đi được trong khoảng thời gian $\dfrac{5T}{4}$ là
\choice
{$2{,}5A$}
{$5A$}
{$A(4+\sqrt{3})$}
{\True $A(4+\sqrt{2})$}
\loigiai{
Ta có
$\dfrac{5T}{4}=T+\dfrac{T}{4}.$
Trong một chu kì, vật đi được quãng đường
$4A.$
Trong thời gian $\dfrac{T}{4}$, quãng đường lớn nhất là
$2A\sin\dfrac{\pi}{4}=A\sqrt{2}.$
Vậy
$S_{\max}=4A+A\sqrt{2}=A(4+\sqrt{2}).$
}
\end{ex}

% ===== Câu 65 =====
% ID: L11C1B1-93-DS
% Mức: VD
\dangbt{Xác định và tính toán các đại lượng vận tốc, gia tốc tại thời điểm xác định}
\begin{ex}
% ID: L11C1B1-93-DS
% Mức: VD
L11C1 Ôn Tập chương dao động – Xác định và tính toán các đại lượng vận tốc, gia tốc tại thời điểm xác định: Từ $x=9\cos(4t+\pi/3)$ cm, hãy lập biểu thức vận tốc và gia tốc.
\choiceTF

\end{ex}

% ===== Câu 66 =====
% ID: L11C1B1-94-TN
% Mức: VDC
\dangbt{quãng đường}
\begin{ex}
% ID: L11C1B1-94-TN
% Mức: VDC
Một chất điểm dao động điều hòa với biên độ $A$ và chu kì $T$. Trong khoảng thời gian $\dfrac{T}{3}$, chất điểm không thể đi được quãng đường bằng
\choice
{$1{,}6A$}
{$1{,}7A$}
{$1{,}5A$}
{\True $1{,}8A$}
\loigiai{
Trong khoảng thời gian $\Delta t=\dfrac{T}{3}$, quãng đường vật đi được thỏa mãn $S_{\min}\le S\le S_{\max}.$

Ta có $S_{\min}=2A\left(1-\cos\dfrac{\pi}{3}\right)=A,S_{\max}=2A\sin\dfrac{\pi}{3}=A\sqrt{3}\approx1,732A.$ Do đó vật không thể đi được quãng đường $1,8A$ trong khoảng thời gian này.
}
\end{ex}

% ===== Câu 67 =====
% ID: L11C1B1-95-TN
% Mức: VDC
\begin{ex}
% ID: L11C1B1-95-TN
% Mức: VDC
Một vật dao động điều hòa với biên độ $A=6\ \text{cm}$ và chu kì $T=1,2\ \text{s}$. Quãng đường lớn nhất mà vật đi được trong khoảng thời gian $2\ \text{s}$ là
\choice
{$34,4\ \text{cm}$}
{\True $42\ \text{cm}$}
{$30\ \text{cm}$}
{$30\sqrt{3}\ \text{cm}$}
\loigiai{
Ta có
$2\ \text{s}=1,2\ \text{s}+0,8\ \text{s}=T+\dfrac{2T}{3}.$
Trong một chu kì, vật đi được
$4A=4\cdot 6=24\ \text{cm}.$
Trong thời gian $\dfrac{2T}{3}$, quãng đường lớn nhất là
$3A=18\ \text{cm}.$
Vậy quãng đường lớn nhất trong $2\ \text{s}$ là
$S_{\max}=24+18=42\ \text{cm}.$
}
\end{ex}

% ===== Câu 68 =====
% ID: L11C1B1-96-TN
% Mức: VDC
\begin{ex}
% ID: L11C1B1-96-TN
% Mức: VDC
Một chất điểm dao động điều hòa dọc theo trục $Ox$. Trong khoảng thời gian $\dfrac{1}{3}\ \text{s}$, vật đi được quãng đường lớn nhất bằng biên độ. Tần số dao động của vật là
\choice
{$2,00\ \text{Hz}$}
{$0,25\ \text{Hz}$}
{$0,75\ \text{Hz}$}
{\True $0,50\ \text{Hz}$}
\loigiai{
Theo đề:
$S_{\max}=A.$
Mà
$S_{\max}=2A\sin\left(\dfrac{\pi\Delta t}{T}\right).$
Suy ra
$2A\sin\left(\dfrac{\pi\Delta t}{T}\right)=A\Rightarrow\sin\left(\dfrac{\pi\Delta t}{T}\right)=\dfrac{1}{2}.$
Lấy nghiệm phù hợp với $\Delta t<\dfrac{T}{2}$: $\dfrac{\pi\Delta t}{T}=\dfrac{\pi}{6}\Rightarrow\dfrac{\Delta t}{T}=\dfrac{1}{6}.$
Với $\Delta t=\dfrac{1}{3}\ \mathrm{s}$: $T=6\Delta t=2\ \mathrm{s}.$ Do đó $f=\dfrac{1}{T}=0,50\ \mathrm{Hz}.$
}
\end{ex}

% ===== Câu 69 =====
% ID: L11C1B1-97-TN
% Mức: VDC
\begin{ex}
% ID: L11C1B1-97-TN
% Mức: VDC
Một vật dao động điều hòa với biên độ $6\ \text{cm}$. Quãng đường nhỏ nhất mà vật đi được trong một giây là $18\ \text{cm}$. Hỏi ở thời điểm kết thúc quãng đường đó thì tốc độ của vật là bao nhiêu?
\choice
{$31,4\ \text{cm/s}$}
{$26,5\ \text{cm/s}$}
{\True $27,2\ \text{cm/s}$}
{$28,1\ \text{cm/s}$}
\loigiai{
Ta có
$A=6\ \text{cm},\qquad S_{\min}=18\ \text{cm}=3A.$
Vì
$3A=2A+A,$
nên thời gian $1\ \text{s}$ ứng với
$\dfrac{T}{2}+\dfrac{T}{3}=\dfrac{5T}{6}.$ Do đó
$\dfrac{5T}{6}=1\Rightarrow T=1,2\ \text{s}.$
Suy ra $\omega=\dfrac{2\pi}{T}=\dfrac{2\pi}{1,2}=\dfrac{5\pi}{3}\ \text{rad/s}$.
Ở thời điểm kết thúc quãng đường nhỏ nhất, vật ở vị trí có
$|x|=\dfrac{A}{2}$.
Tốc độ khi đó:
$v=\omega\sqrt{A^2-x^2}=\omega\sqrt{A^2-\dfrac{A^2}{4}}=\omega\dfrac{A\sqrt{3}}{2}.$
Thay số:
$v=\dfrac{5\pi}{3}\cdot\dfrac{6\sqrt{3}}{2}=5\pi\sqrt{3}\approx 27,2\ \text{cm/s}.$
}
\end{ex}

% ===== Câu 70 =====
% ID: L11C1B1-98-TN
% Mức: VDC
\begin{ex}
% ID: L11C1B1-98-TN
% Mức: VDC
Một chất điểm dao động điều hòa trên trục $Ox$ với chu kì $T$ và biên độ $A$. Vị trí cân bằng của chất điểm trùng với gốc tọa độ. Trong khoảng thời gian $\Delta t$ với $0\lt \Delta t<\dfrac{T}{2}$, quãng đường lớn nhất và nhỏ nhất mà vật có thể đi được lần lượt là $S_{\max}$ và $S_{\min}$. Lựa chọn phương án đúng.
\choice
{$S_{\max}=2A\sin\left(\dfrac{\pi\Delta t}{T}\right);\ S_{\min}=2A\cos\left(\dfrac{\pi\Delta t}{T}\right)$}
{\True $S_{\max}=2A\sin\left(\dfrac{\pi\Delta t}{T}\right);\ S_{\min}=2A-2A\cos\left(\dfrac{\pi\Delta t}{T}\right)$}
{$S_{\max}=2A\sin\left(\dfrac{2\pi\Delta t}{T}\right);\ S_{\min}=2A\cos\left(\dfrac{2\pi\Delta t}{T}\right)$}
{$S_{\max}=2A\sin\left(\dfrac{2\pi\Delta t}{T}\right);\ S_{\min}=2A-2A\cos\left(\dfrac{2\pi\Delta t}{T}\right)$}
\loigiai{
Với $0\lt \Delta t<\dfrac{T}{2}$, đặt
$\Delta\varphi=\omega\Delta t=\dfrac{2\pi\Delta t}{T}.$
Quãng đường lớn nhất ứng với chuyển động đối xứng qua vị trí cân bằng:
 S_{ $\max$ }=2 $A$ \sin\dfrac{\Delta\varphi}{2} $=2$ A $\sin\left(\dfrac{\pi\Delta t}{T}\right)$. 
Quãng đường nhỏ nhất ứng với chuyển động gần biên:
 S_{ $\min$ }=2 $A$ \left(1-\cos\dfrac{\Delta\varphi}{2}\right) $=$ 2\,\text{A} $-2$ A $\cos\left(\dfrac{\pi\Delta t}{T}\right)$.
}
\end{ex}

% ===== Câu 71 =====
% ID: L11C1B1-99-TN
% Mức: VDC
\begin{ex}
% ID: L11C1B1-99-TN
% Mức: VDC
Một vật dao động điều hòa dọc theo trục $Ox$, quanh vị trí cân bằng $O$ với biên độ $A$. Trong khoảng thời gian $1\ \text{s}$, quãng đường nhỏ nhất mà vật có thể đi được là $A$. Chu kì dao động điều hòa là
\choice
{$5\ \text{s}$}
{\True $3\ \text{s}$}
{$4\ \text{s}$}
{$2,5\ \text{s}$}
\loigiai{
Ta có
$S_{\min}=2A\left[1-\cos\left(\dfrac{\pi\Delta t}{T}\right)\right].$
Theo đề:
$2A\left[1-\cos\left(\dfrac{\pi}{T}\right)\right]=A.$
Suy ra
$1-\cos\left(\dfrac{\pi}{T}\right)=\dfrac{1}{2}\Rightarrow\cos\left(\dfrac{\pi}{T}\right)=\dfrac{1}{2}.$
Do đó
$\dfrac{\pi}{T}=\dfrac{\pi}{3}\Rightarrow T=3\ \mathrm{s}.$
}
\end{ex}

% ===== Câu 72 =====
% ID: L11C1B1-100-TN
% Mức: VDC
\begin{ex}
% ID: L11C1B1-100-TN
% Mức: VDC
Một vật dao động điều hòa với biên độ $4\ \text{cm}$. Quãng đường nhỏ nhất mà vật đi được trong $1\ \text{s}$ là $20\ \text{cm}$. Gia tốc lớn nhất của vật là bao nhiêu? Lấy $\pi^2=10$.
\choice
{$4,82\ \text{m/s}^2$}
{$248,42\ \text{cm/s}^2$}
{$3,96\ \text{m/s}^2$}
{\True $284,44\ \text{cm/s}^2$}
\loigiai{
Ta có
$A=4\ \text{cm},\qquad S_{\min}=20\ \text{cm}=5A.$
Vì trong một chu kì vật đi được $4A$, ta viết
$5A=4A+A.$
Do đó
$1\ \text{s}=T+\dfrac{T}{3}=\dfrac{4T}{3}.$
Suy ra
$T=0,75\ \text{s}.$
Tần số góc:
 $\omega=\dfrac{2\pi}{T}$ = $\dfrac{2\pi}{0,75}$ = $\dfrac{8\pi}{3}\\text{rad/s}$. 
Gia tốc cực đại:
 a_{ $\max$ }= $\omega^2A$ = $\left(\dfrac{8\pi}{3}\right)^2\cdot$ 4
= $\dfrac{256\pi^2}{9}$. 
Lấy $\pi^2=10$:
$a_{\max}=\dfrac{2560}{9}\approx 284,44\ \text{cm/s}^2.$
}
\end{ex}

% ===== Câu 73 =====
% ID: L11C1B1-101-TN
% Mức: TH
\dangbt{Thời điểm vật qua x1}
\begin{ex}
% ID: L11C1B1-101-TN
% Mức: TH
Một chất điểm dao động điều hòa với chu kì $T$. Khoảng thời gian trong một chu kì để vật có tọa độ âm là
\choice
{$\dfrac{T}{3}$}
{$\dfrac{2T}{3}$}
{$\dfrac{T}{6}$}
{\True $\dfrac{T}{2}$}
\loigiai{
Trong dao động điều hòa quanh vị trí cân bằng, trong một chu kì vật ở phía âm của trục tọa độ trong nửa chu kì.
Do đó
$\Delta t=\dfrac{T}{2}.$
}
\end{ex}

% ===== Câu 74 =====
% ID: L11C1B1-102-TN
% Mức: TH
\begin{ex}
% ID: L11C1B1-102-TN
% Mức: TH
Một chất điểm dao động điều hòa với chu kì $T$ trên trục $Ox$, với $O$ là vị trí cân bằng. Thời gian ngắn nhất vật đi từ điểm có tọa độ $x=0$ đến điểm có tọa độ $x=\dfrac{A}{2}$ là
\choice
{$\dfrac{T}{24}$}
{$\dfrac{T}{16}$}
{$\dfrac{T}{6}$}
{\True $\dfrac{T}{12}$}
\loigiai{
Từ $x=0$ đến $x=\dfrac{A}{2}$, ta có
$\sin\Delta\varphi=\dfrac{x}{A}=\dfrac{1}{2}.$
Do đó
$\Delta\varphi=\dfrac{\pi}{6}.$
Suy ra
 $\Delta$ t= $\dfrac{\Delta\varphi}{2\pi}T$ = $\dfrac{\pi/6}{2\pi}T$ = $\dfrac{T}{12}$.
}
\end{ex}

% ===== Câu 75 =====
% ID: L11C1B1-103-TN
% Mức: TH
\begin{ex}
% ID: L11C1B1-103-TN
% Mức: TH
Vật dao động điều hòa, thời gian ngắn nhất vật đi từ vị trí cân bằng đến vị trí có li độ cực đại là $0,1\ \text{s}$. Chu kì dao động của vật là
\choice
{$0,05\ \text{s}$}
{$0,1\ \text{s}$}
{$0,2\ \text{s}$}
{\True $0,4\ \text{s}$}
\loigiai{
Thời gian ngắn nhất vật đi từ vị trí cân bằng đến vị trí biên là $\dfrac{T}{4}$.

Theo đề: $\dfrac{T}{4}=0,1\Rightarrow T=0,4\ \mathrm{s}$.
}
\end{ex}

% ===== Câu 76 =====
% ID: L11C1B1-104-TN
% Mức: TH
\begin{ex}
% ID: L11C1B1-104-TN
% Mức: TH
Một chất điểm dao động điều hòa với chu kì $T$ trên trục $Ox$, với $O$ là vị trí cân bằng. Thời gian ngắn nhất vật đi từ tọa độ $x=0$ đến tọa độ $x=\dfrac{A}{\sqrt{2}}$ là
\choice
{\True $\dfrac{T}{8}$}
{$\dfrac{T}{16}$}
{$\dfrac{T}{6}$}
{$\dfrac{T}{12}$}
\loigiai{
Từ $x=0$ đến $x=\dfrac{A}{\sqrt{2}}$, ta có $\sin\Delta\varphi=\dfrac{1}{\sqrt{2}}.$

Suy ra $\Delta\varphi=\dfrac{\pi}{4}.$

Vậy $\Delta t=\dfrac{\Delta\varphi}{2\pi}T=\dfrac{\pi/4}{2\pi}T=\dfrac{T}{8}$.
}
\end{ex}

% ===== Câu 77 =====
% ID: L11C1B1-105-TN
% Mức: VD
\begin{ex}
% ID: L11C1B1-105-TN
% Mức: VD
Khoảng thời gian ngắn nhất để vật dao động điều hòa chuyển động từ li độ $x_1=-\dfrac{A}{2}$ đến $x_2=0,5A\sqrt{3}$ là
\choice
{\True $\dfrac{T}{4}$}
{$\dfrac{T}{3}$}
{$\dfrac{T}{2}$}
{$\dfrac{T}{6}$}
\loigiai{
Ta có
$x_1=-\dfrac{A}{2},\qquad x_2=\dfrac{A\sqrt{3}}{2}.$
Góc quét ngắn nhất giữa hai vị trí này trên vòng tròn lượng giác là
$\Delta\varphi=\dfrac{\pi}{2}.$
Suy ra
 $\Delta$ t= $\dfrac{\Delta\varphi}{2\pi}T$ = $\dfrac{\pi/2}{2\pi}T$ = $\dfrac{T}{4}$.
}
\end{ex}

% ===== Câu 78 =====
% ID: L11C1B1-106-TN
% Mức: VD
\begin{ex}
% ID: L11C1B1-106-TN
% Mức: VD
Một chất điểm dao động điều hòa với chu kì $T$. Khoảng thời gian trong một chu kì để vật cách vị trí cân bằng một khoảng nhỏ hơn $0,5\sqrt{3}$ biên độ là
\choice
{$\dfrac{T}{3}$}
{\True $\dfrac{2T}{3}$}
{$\dfrac{T}{6}$}
{$\dfrac{T}{2}$}
\loigiai{
Ta có
$0,5\sqrt{3}A=\dfrac{A\sqrt{3}}{2}.$
Điều kiện:
$|x|\lt \dfrac{A\sqrt{3}}{2}.$
Trong một chu kì, tổng góc quét thỏa điều kiện này là
$\Delta\varphi=\dfrac{4\pi}{3}.$
Vậy
 $\Delta$ t= $\dfrac{4\pi/3}{2\pi}T$ = $\dfrac{2T}{3}$.
}
\end{ex}

% ===== Câu 79 =====
% ID: L11C1B1-107-TN
% Mức: VD
\begin{ex}
% ID: L11C1B1-107-TN
% Mức: VD
Khoảng thời gian ngắn nhất mà một vật dao động điều hòa với chu kì $T$, biên độ $A$ thực hiện khi di chuyển giữa hai vị trí có li độ $x_1=\dfrac{A}{2}$ và $x_2=0,5A\sqrt{3}$ là
\choice
{$\dfrac{T}{6}$}
{$\dfrac{T}{8}$}
{$0,5T(\sqrt{3}-1)$}
{\True $\dfrac{T}{12}$}
\loigiai{
Ta có
$x_1=\dfrac{A}{2},\qquad x_2=\dfrac{A\sqrt{3}}{2}.$
Góc tương ứng với $x_1$ là $\dfrac{\pi}{3}$, với $x_2$ là $\dfrac{\pi}{6}$.
Góc quét ngắn nhất:
 $\Delta\varphi=\dfrac{\pi}{3}-\dfrac{\pi}{6}$ = $\dfrac{\pi}{6}$. 
Vậy
 $\Delta$ t= $\dfrac{\Delta\varphi}{2\pi}T$ = $\dfrac{\pi/6}{2\pi}T$ = $\dfrac{T}{12}$.
}
\end{ex}

% ===== Câu 80 =====
% ID: L11C1B1-108-TN
% Mức: VDC
\begin{ex}
% ID: L11C1B1-108-TN
% Mức: VDC
Một chất điểm đang dao động điều hòa trên một đoạn thẳng. Trên đoạn thẳng đó có bảy điểm theo đúng thứ tự $M_1,M_2,M_3,M_4,M_5,M_6,M_7$, với $M_4$ là vị trí cân bằng. Biết cứ $0,05\ \text{s}$ thì chất điểm lại đi qua các điểm $M_1,M_2,M_3,M_4,M_5,M_6,M_7$; tốc độ tại $M_1$ và $M_7$ bằng $0$. Tốc độ của nó lúc đi qua điểm $M_4$ là $20\pi\ \text{cm/s}$. Biên độ $A$ bằng
\choice
{$4\ \text{cm}$}
{\True $6\ \text{cm}$}
{$4\sqrt{2}\ \text{cm}$}
{$4\sqrt{3}\ \text{cm}$}
\loigiai{
Theo bài trước: $T = 0,6\ \mathrm{s}$.

Do đó $\omega = \dfrac{2\pi}{T} = \dfrac{2\pi}{0,6} = \dfrac{10\pi}{3}\ \mathrm{rad/s}$.

Tại vị trí cân bằng $M_4$, tốc độ của vật là cực đại: $v_{\max} = \omega A$.

Theo đề: $20\pi = \dfrac{10\pi}{3} A \Rightarrow A = 6\ \mathrm{cm}$.
}
\end{ex}

% ===== Câu 81 =====
% ID: L11C1B1-109-TN
% Mức: VDC
\begin{ex}
% ID: L11C1B1-109-TN
% Mức: VDC
Một chất điểm đang dao động điều hòa trên một đoạn thẳng. Trên đoạn thẳng đó có năm điểm theo đúng thứ tự $M$, $N$, $O$, $P$ và $Q$, với $O$ là vị trí cân bằng. Biết cứ $0,05\ \text{s}$ thì chất điểm lại đi qua các điểm $M$, $N$, $O$, $P$ và $Q$; tốc độ tại $M$ và $Q$ bằng $0$. Tốc độ của nó lúc đi qua các điểm $N$, $P$ là $20\pi\ \text{cm/s}$. Biên độ $A$ bằng
\choice
{$4\ \text{cm}$}
{$6\ \text{cm}$}
{\True $4\sqrt{2}\ \text{cm}$}
{$4\sqrt{3}\ \text{cm}$}
\loigiai{
Theo bài trước:
$T=0,4\ \mathrm{s}.$
Do đó
$\omega=\dfrac{2\pi}{T}=\dfrac{2\pi}{0,4}=5\pi\ \mathrm{rad/s}.$
Từ biên $M$ đến biên $Q$ chia thành $4$ khoảng thời gian bằng nhau nên các pha cách nhau $\dfrac{\pi}{4}$. Tại $N$ hoặc $P$:
$|x|=\dfrac{A}{\sqrt{2}}.$
Suy ra tốc độ tại $N$ hoặc $P$:
$v=\omega\sqrt{A^2-x^2}=\omega\sqrt{A^2-\dfrac{A^2}{2}}=\dfrac{\omega A}{\sqrt{2}}.$
Theo đề:
$\dfrac{5\pi A}{\sqrt{2}}=20\pi\Rightarrow A=4\sqrt{2}\ \mathrm{cm}.$
}
\end{ex}

% ===== Câu 82 =====
% ID: L11C1B1-110-TN
% Mức: VDC
\begin{ex}
% ID: L11C1B1-110-TN
% Mức: VDC
Một dao động điều hòa có chu kì dao động là $T$ và biên độ là $A$. Tại thời điểm ban đầu vật có li độ $x_0\gt 0$. Thời gian ngắn nhất để vật đi từ vị trí ban đầu về vị trí cân bằng chỉ bằng một nửa thời gian ngắn nhất để vật đi từ vị trí ban đầu về vị trí biên $x=+A$. Chọn phương án đúng.
\choice
{$x_0=0{,}25A$}
{$x_0=0,5A\sqrt{3}$}
{$x_0=0,5A\sqrt{2}$}
{\True $x_0=0{,}5A$}
\loigiai{
Đặt
$x_0=A\cos\alpha,\qquad 0\lt \alpha<\dfrac{\pi}{2}.$
Thời gian đến biên dương:
$t_1=\dfrac{\alpha}{\omega}.$
Thời gian đến vị trí cân bằng:
$t_2=\dfrac{\frac{\pi}{2}-\alpha}{\omega}.$
Theo đề:
$t_2=\dfrac{1}{2}t_1.$
Suy ra
 $\dfrac{\pi}{2}-\alpha=\dfrac{\alpha}{2}\Rightarrow\alpha=\dfrac{\pi}{3}$. 
Vậy
 x_0= $A$ \cos\dfrac{\pi}{3} $=$ \dfrac{A}{2} $=$ 0,5\,\text{A}.
}
\end{ex}

% ===== Câu 83 =====
% ID: L11C1B1-111-TN
% Mức: VDC
\begin{ex}
% ID: L11C1B1-111-TN
% Mức: VDC
Một dao động điều hòa có chu kì dao động là $T$ và biên độ là $A$. Tại thời điểm ban đầu vật có li độ $x_0\gt 0$. Thời gian ngắn nhất để vật đi từ vị trí ban đầu về vị trí cân bằng gấp đôi thời gian ngắn nhất để vật đi từ vị trí ban đầu về vị trí biên $x=+A$. Chọn phương án đúng.
\choice
{$x_0=0{,}25A$}
{\True $x_0=0,5A\sqrt{3}$}
{$x_0=0,5A\sqrt{2}$}
{$x_0=0{,}5A$}
\loigiai{
Đặt
$x_0=A\cos\alpha,\qquad 0\lt \alpha<\dfrac{\pi}{2}.$
Thời gian đến biên dương:
$t_1=\dfrac{\alpha}{\omega}.$
Thời gian đến vị trí cân bằng:
$t_2=\dfrac{\frac{\pi}{2}-\alpha}{\omega}.$
Theo đề:
$t_2=2t_1.$
Suy ra
 $\dfrac{\pi}{2}-\alpha=2\alpha\Rightarrow\alpha=\dfrac{\pi}{6}$. 
Vậy
 x_0= $A$ \cos\dfrac{\pi}{6} $=$ \dfrac{A\sqrt{3}}{2} $=0,5$ A $\sqrt{3}$.
}
\end{ex}

% ===== Câu 84 =====
% ID: L11C1B1-112-TN
% Mức: VDC
\begin{ex}
% ID: L11C1B1-112-TN
% Mức: VDC
Một dao động điều hòa có chu kì dao động là $T$ và biên độ là $A$. Tại thời điểm ban đầu vật có li độ $x_0\gt 0$. Thời gian ngắn nhất để vật đi từ vị trí ban đầu về vị trí cân bằng cũng bằng thời gian ngắn nhất để vật đi từ vị trí ban đầu về vị trí biên $x=+A$. Chọn phương án đúng.
\choice
{$x_0=0{,}25A$}
{$x_0=0,5A\sqrt{3}$}
{\True $x_0=0,5A\sqrt{2}$}
{$x_0=0{,}5A$}
\loigiai{
Đặt
$x_0=A\cos\alpha,\qquad 0\lt \alpha<\dfrac{\pi}{2}.$
Thời gian đến biên dương:
$t_1=\dfrac{\alpha}{\omega}.$
Thời gian đến vị trí cân bằng:
$t_2=\dfrac{\frac{\pi}{2}-\alpha}{\omega}.$
Theo đề:
$t_1=t_2.$
Suy ra
 $\alpha=\dfrac{\pi}{2}-\alpha\Rightarrow\alpha=\dfrac{\pi}{4}$. 
Do đó
 x_0= $A$ \cos\dfrac{\pi}{4} $=$ \dfrac{A\sqrt{2}}{2} $=0,5$ A $\sqrt{2}$.
}
\end{ex}

% ===== Câu 85 =====
% ID: L11C1B1-113-TN
% Mức: VDC
\begin{ex}
% ID: L11C1B1-113-TN
% Mức: VDC
Một chất điểm đang dao động điều hòa trên một đoạn thẳng xung quanh vị trí cân bằng $O$. Gọi $M$, $N$ là hai điểm trên đường thẳng cùng cách đều $O$. Biết cứ $0,05\ \text{s}$ thì chất điểm lại đi qua các điểm $M$, $O$, $N$ và tốc độ tại $M$ và $N$ khác $0$. Chu kì dao động bằng
\choice
{\True $0,3\ \text{s}$}
{$0,4\ \text{s}$}
{$0,2\ \text{s}$}
{$0,1\ \text{s}$}
\loigiai{
Vật lần lượt đi qua $M$, $O$, $N$ trong các khoảng thời gian bằng nhau $0,05\ \mathrm{s}$. Do $M$ và $N$ đối xứng qua $O$, các vị trí này tương ứng với các pha cách nhau $\Delta\varphi = \dfrac{\pi}{3}$. Vậy $0,05 = \dfrac{\Delta\varphi}{2\pi}T = \dfrac{\pi/3}{2\pi}T = \dfrac{T}{6}$. Suy ra $T = 0,3\ \mathrm{s}$.
}
\end{ex}

% ===== Câu 86 =====
% ID: L11C1B1-114-TN
% Mức: VDC
\begin{ex}
% ID: L11C1B1-114-TN
% Mức: VDC
Một chất điểm đang dao động điều hòa trên một đoạn thẳng. Trên đoạn thẳng đó có bảy điểm theo đúng thứ tự $M_1,M_2,M_3,M_4,M_5,M_6,M_7$, với $M_4$ là vị trí cân bằng. Biết cứ $0,05\ \text{s}$ thì chất điểm lại đi qua các điểm $M_1,M_2,M_3,M_4,M_5,M_6,M_7$; tốc độ tại $M_1$ và $M_7$ bằng $0$. Chu kì bằng
\choice
{$0,3\ \text{s}$}
{$0,4\ \text{s}$}
{$0,2\ \text{s}$}
{\True $0,6\ \text{s}$}
\loigiai{
Vì tốc độ tại $M_1$ và $M_7$ bằng $0$ nên $M_1$, $M_7$ là hai vị trí biên.
Từ $M_1$ đến $M_7$ là nửa chu kì và gồm $6$ khoảng thời gian bằng nhau:
$\dfrac{T}{2}=6\cdot 0,05=0,30\ \text{s}.$
Suy ra
$T=0,60\ \text{s}.$
}
\end{ex}

% ===== Câu 87 =====
% ID: L11C1B1-115-TN
% Mức: TH
\dangbt{Viết phương trình dao động điều hòa.}
\begin{ex}
% ID: L11C1B1-115-TN
% Mức: TH
Hình chiếu của một chất điểm chuyển động tròn đều lên một đường kính quỹ đạo có chuyển động là dao động điều hòa. Phát biểu nào sau đây sai?
\choice
{Tần số góc của dao động điều hòa bằng tốc độ góc của chuyển động tròn đều}
{Biên độ của dao động điều hòa bằng bán kính của chuyển động tròn đều}
{\True Lực kéo về trong dao động điều hòa có độ lớn bằng độ lớn lực hướng tâm trong chuyển động tròn đều}
{Tốc độ cực đại của dao động điều hòa bằng tốc độ dài của chuyển động tròn đều}
\loigiai{
Hình chiếu của chuyển động tròn đều lên một đường kính là dao động điều hòa.

Ta có
$A=R,\qquad \omega_{\text{dao động}}=\omega_{\text{tròn đều}},\qquad v_{\max}=\omega R.$
Tuy nhiên, lực kéo về trong dao động điều hòa có độ lớn
$F=m\omega^2|x|,$
còn lực hướng tâm trong chuyển động tròn đều có độ lớn
$F_{\text{ht}}=m\omega^2R.$
Hai lực này không luôn bằng nhau. Vậy phát biểu sai là phương án C.
}
\end{ex}

% ===== Câu 88 =====
% ID: L11C1B1-116-TN
% Mức: TH
\begin{ex}
% ID: L11C1B1-116-TN
% Mức: TH
Một vật dao động điều hòa với biên độ $A$ và tốc độ cực đại $v_{\max}$. Tần số góc của vật dao động là
\choice
{\True $\dfrac{v_{\max}}{A}$}
{$\dfrac{v_{\max}}{\pi A}$}
{$\dfrac{v_{\max}}{2\pi A}$}
{$\dfrac{v_{\max}}{2A}$}
\loigiai{
Ta có
$v_{\max}=\omega A.$
Suy ra
$\omega=\dfrac{v_{\max}}{A}.$
}
\end{ex}

% ===== Câu 89 =====
% ID: L11C1B1-117-TN
% Mức: TH
\begin{ex}
% ID: L11C1B1-117-TN
% Mức: TH
Khi nói về một vật đang dao động điều hòa, phát biểu nào sau đây đúng?
\choice
{Vectơ gia tốc của vật đổi chiều khi vật có li độ cực đại}
{\True Vectơ vận tốc và vectơ gia tốc của vật cùng chiều nhau khi vật chuyển động về phía vị trí cân bằng}
{Vectơ gia tốc của vật luôn hướng ra xa vị trí cân bằng}
{Vectơ vận tốc và vectơ gia tốc của vật cùng chiều nhau khi vật chuyển động ra xa vị trí cân bằng}
\loigiai{
Trong dao động điều hòa:
$a=-\omega^2x.$
Gia tốc luôn hướng về vị trí cân bằng. Khi vật chuyển động về vị trí cân bằng, vận tốc cũng hướng về vị trí cân bằng. Do đó vectơ vận tốc và vectơ gia tốc cùng chiều nhau.
}
\end{ex}

% ===== Câu 90 =====
% ID: L11C1B1-118-TN
% Mức: TH
\begin{ex}
% ID: L11C1B1-118-TN
% Mức: TH
Khi một vật dao động điều hòa, chuyển động của vật từ vị trí biên về vị trí cân bằng là chuyển động
\choice
{nhanh dần đều}
{chậm dần đều}
{\True nhanh dần}
{chậm dần}
\loigiai{
Khi vật đi từ vị trí biên về vị trí cân bằng, tốc độ tăng từ $0$ đến $v_{\max}$ nên vật chuyển động nhanh dần. Tuy nhiên gia tốc không không đổi nên không phải nhanh dần đều.
}
\end{ex}

% ===== Câu 91 =====
% ID: L11C1B1-119-TN
% Mức: VD
\begin{ex}
% ID: L11C1B1-119-TN
% Mức: VD
Một vật dao động điều hòa với tần số góc $5\ \text{rad/s}$. Khi vật đi qua li độ $5\ \text{cm}$ thì nó có tốc độ là $25\ \text{cm/s}$. Biên độ dao động của vật là
\choice
{$5,24\ \text{cm}$}
{\True $5\sqrt{2}\ \text{cm}$}
{$5\sqrt{3}\ \text{cm}$}
{$10\ \text{cm}$}
\loigiai{
Ta có
$v^2=\omega^2\left(A^2-x^2\right).$
Thay số:
$25^2=5^2\left(A^2-5^2\right).$
Suy ra
$25=A^2-25\Rightarrow A^2=50\Rightarrow A=5\sqrt{2}\ \mathrm{cm}$.
}
\end{ex}

% ===== Câu 92 =====
% ID: L11C1B1-120-TN
% Mức: VD
\begin{ex}
% ID: L11C1B1-120-TN
% Mức: VD
Một vật nhỏ dao động điều hòa theo phương trình
$x=A\cos 4\pi t$
trong đó $t$ tính bằng giây. Tính từ $t=0$, khoảng thời gian ngắn nhất để gia tốc của vật có độ lớn bằng một nửa độ lớn gia tốc cực đại là
\choice
{\True $0,083\ \text{s}$}
{$0,125\ \text{s}$}
{$0,104\ \text{s}$}
{$0,167\ \text{s}$}
\loigiai{
Trong dao động điều hòa:
$a = -\omega^2 x.$ Do đó
$|a| = \dfrac{1}{2}a_{\max} \Rightarrow |x| = \dfrac{A}{2}.$
Vì $x = A\cos 4\pi t,$
nên
$|\cos 4\pi t| = \dfrac{1}{2}.$
Tính từ $t = 0$, thời điểm sớm nhất là
$4\pi t = \dfrac{\pi}{3} \Rightarrow t = \dfrac{1}{12}\ \text{s} \approx 0,083\ \text{s}.$
}
\end{ex}

% ===== Câu 93 =====
% ID: L11C1B1-121-TN
% Mức: VD
\begin{ex}
% ID: L11C1B1-121-TN
% Mức: VD
Một vật nhỏ dao động điều hòa với biên độ $4\ \text{cm}$ và chu kì $2\ \text{s}$. Quãng đường vật đi được trong $4\ \text{s}$ là
\choice
{$8\ \text{cm}$}
{$16\ \text{cm}$}
{$64\ \text{cm}$}
{\True $32\ \text{cm}$}
\loigiai{
Trong một chu kì, vật đi được quãng đường
$S_T=4A.$
Vì
$4\ \text{s}=2T,$
nên
$S=2\cdot 4A=8A=8\cdot 4=32\ \text{cm}.$
}
\end{ex}

% ===== Câu 94 =====
% ID: L11C1B1-122-TN
% Mức: VD
\begin{ex}
% ID: L11C1B1-122-TN
% Mức: VD
Một vật dao động điều hòa có chu kì $2\ \text{s}$, biên độ $10\ \text{cm}$. Khi vật cách vị trí cân bằng $6\ \text{cm}$, tốc độ của nó bằng
\choice
{$18,84\ \text{cm/s}$}
{$20,08\ \text{cm/s}$}
{\True $25,13\ \text{cm/s}$}
{$12,56\ \text{cm/s}$}
\loigiai{
Ta có
$\omega=\dfrac{2\pi}{T}=\dfrac{2\pi}{2}=\pi\ \mathrm{rad/s}.$
Tốc độ tại li độ $x$:
$v=\omega\sqrt{A^2-x^2}.$
Thay số:
$v=\pi\sqrt{10^2-6^2}=\pi\sqrt{64}=8\pi\approx 25,13\ \mathrm{cm/s}.$
}
\end{ex}

% ===== Câu 95 =====
% ID: L11C1B1-123-TN
% Mức: VDC
\begin{ex}
% ID: L11C1B1-123-TN
% Mức: VDC
Một thấu kính hội tụ có tiêu cự $15\ \text{cm}$. $M$ là một điểm nằm trên trục chính của thấu kính, $P$ là một chất điểm dao động điều hòa quanh vị trí cân bằng trùng với $M$. Gọi $P'$ là ảnh của $P$ qua thấu kính. Khi $P$ dao động theo phương vuông góc với trục chính, biên độ $5\ \text{cm}$ thì $P'$ là ảnh ảo dao động với biên độ $10\ \text{cm}$. Nếu $P$ dao động dọc theo trục chính với tần số $5\ \text{Hz}$, biên độ $2,5\ \text{cm}$ thì $P'$ có tốc độ trung bình trong khoảng thời gian $0,2\ \text{s}$ bằng
\choice
{$1,5\ \text{m/s}$}
{$1,25\ \text{m/s}$}
{\True $2,25\ \text{m/s}$}
{$1,0\ \text{m/s}$}
\loigiai{
Độ phóng đại ảnh khi vật dao động vuông góc với trục chính là
$|k|=\dfrac{10}{5}=2.$
Vì ảnh là ảnh ảo nên
$d'=-2d.$
Theo công thức thấu kính:
$\dfrac{1}{f}=\dfrac{1}{d}+\dfrac{1}{d'}.$
Thay $f=15\ \text{cm}$ và $d'=-2d$:
 $\dfrac{1}{15}=\dfrac{1}{d}-\dfrac{1}{2d}$ = $\dfrac{1}{2d}$. 
Suy ra
$d=7,5\ \text{cm}.$

Khi vật dao động dọc theo trục chính với biên độ $2,5\ \text{cm}$:
$d_{\min}=7,5-2,5=5\ \text{cm},$ d_{ $\max$ }=7,5+2,5=10 $\\text{cm}$. $Vị trí ảnh:$ d'= $\dfrac{fd}{d-f}$. $Với$ d=5 $\\text{cm}:$ d'_1= $\dfrac{15\cdot 5}{5-15}=-7,5\\text{cm}$. $Với$ d=10 $\\text{cm}:$ d'_2= $\dfrac{15\cdot 10}{10-15}=-30\\text{cm}$. $Trong một chu kì, ảnh đi được quãng đường$ S=2 $\left|-30-(-7,5)\right|=45\\text{cm}=0,45\\text{m}$. $Vì$ T= $\dfrac{1}{f}=\dfrac{1}{5}=0,2\\text{s},nên tốc độ trung bình là$ v_{ $\text{tb}$ }= $\dfrac{0,45}{0,2}=2,25\\text{m/s}$.
}
\end{ex}

% ===== Câu 96 =====
% ID: L11C1B1-124-TN
% Mức: VDC
\begin{ex}
% ID: L11C1B1-124-TN
% Mức: VDC
Con lắc lò xo gồm một vật nhỏ có khối lượng $250\ \text{g}$ và lò xo nhẹ có độ cứng $100\ \text{N/m}$ dao động điều hòa dọc theo trục $Ox$ với biên độ $4\ \text{cm}$. Khoảng thời gian ngắn nhất để vận tốc của vật có giá trị từ $-40\ \text{cm/s}$ đến $40\sqrt{3}\ \text{cm/s}$ là
\choice
{\True $\dfrac{\pi}{40}\ \text{s}$}
{$\dfrac{\pi}{120}\ \text{s}$}
{$\dfrac{\pi}{20}\ \text{s}$}
{$\dfrac{\pi}{60}\ \text{s}$}
\loigiai{
Ta có
 $\omega=\sqrt{\dfrac{k}{m}}=\sqrt{\dfrac{100}{0,25}}=20\\text{rad/s}$. 
Lại có
$v_{\max}=\omega A=20\cdot 4=80\ \text{cm/s}.$

Khoảng biến thiên vận tốc từ
$v=-40=-\dfrac{v_{\max}}{2}$
đến
$v=40\sqrt{3}=\dfrac{\sqrt{3}}{2}v_{\max}.$
Góc quét nhỏ nhất tương ứng là
$\Delta\varphi=\dfrac{\pi}{2}.$
Vậy
 $\Delta$ t= $\dfrac{\Delta\varphi}{\omega}$ = $\dfrac{\pi/2}{20}$ = $\dfrac{\pi}{40}\\text{s}$.
}
\end{ex}

% ===== Câu 97 =====
% ID: L11C1B1-125-TN
% Mức: VDC
\begin{ex}
% ID: L11C1B1-125-TN
% Mức: VDC
Hai vật dao động điều hòa dọc theo các trục song song với nhau. Phương trình dao động của các vật lần lượt là
$x_1=A_1 \cos\omega t \text{ (cm)}, \qquad x_2=A_2 \sin\omega t \text{ (cm)}.$

Biết
$64x_1^2+36x_2^2=48^2 \text{ (cm}^2\text{)}.$

Tại thời điểm $t$, vật thứ nhất đi qua vị trí có li độ $x_1=3$ cm với vận tốc $v_1=-18$ cm/s. Khi đó vật thứ hai có tốc độ bằng
\choice
{$24\sqrt{3}\ \text{cm/s}$}
{$24\ \text{cm/s}$}
{$8\ \text{cm/s}$}
{\True $8\sqrt{3}\ \text{cm/s}$}
\loigiai{
Từ
$64x_1^2+36x_2^2=48^2$
suy ra
$64A_1^2=48^2$, $36A_2^2=48^2$.
Do đó
$A_1=6\ \mathrm{cm}$, $A_2=8\ \mathrm{cm}$.
Vì $x_1=A_1\cos\omega t=3$ suy ra $\cos\omega t=\dfrac{1}{2}$.
Mặt khác $v_1=-A_1\omega\sin\omega t=-18$.
Do $v_1\lt 0$ nên $\sin\omega t=\dfrac{\sqrt{3}}{2}$.
Suy ra $18=6\omega\cdot\dfrac{\sqrt{3}}{2}$, do đó $\omega=2\sqrt{3}\ \mathrm{rad/s}$.

Vận tốc của vật thứ hai: $v_2=A_2\omega\cos\omega t$.
Do đó $|v_2|=8\cdot 2\sqrt{3}\cdot\dfrac{1}{2}=8\sqrt{3}\ \mathrm{cm/s}$.
}
\end{ex}

% ===== Câu 98 =====
% ID: L11C1B1-126-TN
% Mức: TH
\dangbt{vận tốc trung bình và tốc độ trung bình}
\begin{ex}
% ID: L11C1B1-126-TN
% Mức: TH
Một vật dao động điều hòa với biên độ $10\,\mathrm{cm}$ và tần số $2\,\mathrm{Hz}$. Tại thời điểm $t=0$, vật chuyển động theo chiều dương và đến thời điểm $t=2\,\mathrm{s}$ vật có gia tốc $80\pi^2\sqrt{2}\,\mathrm{cm/s^2}$. Tốc độ trung bình của vật từ thời điểm $t_1=0,0625\,\mathrm{s}$ đến thời điểm $t_2=0,1875\,\mathrm{s}$ gần nhất với giá trị nào sau đây?
\choice
{$99\,\mathrm{cm/s}$}
{$40\,\mathrm{cm/s}$}
{\True $80\,\mathrm{cm/s}$}
{$65\,\mathrm{cm/s}$}
\loigiai{
Dao động điều hòa có biên độ $A = 10\,\mathrm{cm}$ và tần số $f = 2\,\mathrm{Hz}$. Chu kỳ dao động là $T = \frac{1}{f} = \frac{1}{2} = 0,5\,\mathrm{s}$. Tần số góc là $\omega = 2\pi f = 4\pi\,\mathrm{rad/s}$.
Phương trình dao động có dạng $x = A\cos(\omega t + \phi)$.
Tại thời điểm $t=0$, vật chuyển động theo chiều dương nên vận tốc $v \gt  0$. Vận tốc là $v = -\omega A\sin(\omega t + \phi)$. Tại $t=0$, $v = -\omega A\sin(\phi) \gt  0$, suy ra $\sin(\phi) \lt  0$.
Gia tốc là $a = -\omega^2 x = -\omega^2 A\cos(\omega t + \phi)$.
Tại thời điểm $t=2\,\mathrm{s}$, gia tốc $a = 80\pi^2\sqrt{2}\,\mathrm{cm/s^2}$.
$a = -(4\pi)^2 \cdot 10 \cos(4\pi \cdot 2 + \phi) = -160\pi^2 \cos(8\pi + \phi) = -160\pi^2 \cos(\phi)$.
Ta có $-160\pi^2 \cos(\phi) = 80\pi^2\sqrt{2}$, suy ra $\cos(\phi) = -\frac{\sqrt{2}}{2}$.
Kết hợp với $\sin(\phi) \lt  0$, ta chọn $\phi = -\frac{3\pi}{4}$.
Phương trình dao động là $x = 10\cos(4\pi t - \frac{3\pi}{4})\,\mathrm{cm}$.
Vận tốc là $v = -40\pi\sin(4\pi t - \frac{3\pi}{4})\,\mathrm{cm/s}$.
Tại thời điểm $t_1 = 0,0625\,\mathrm{s} = \frac{1}{16}\,\mathrm{s}$:
$4\pi t_1 - \frac{3\pi}{4} = 4\pi \cdot \frac{1}{16} - \frac{3\pi}{4} = \frac{\pi}{4} - \frac{3\pi}{4} = -\frac{2\pi}{4} = -\frac{\pi}{2}$.
$x_1 = 10\cos(-\frac{\pi}{2}) = 0\,\mathrm{cm}$.
$v_1 = -40\pi\sin(-\frac{\pi}{2}) = -40\pi(-1) = 40\pi\,\mathrm{cm/s}$.
Tại thời điểm $t_2 = 0,1875\,\mathrm{s} = \frac{3}{16}\,\mathrm{s}$:
$4\pi t_2 - \frac{3\pi}{4} = 4\pi \cdot \frac{3}{16} - \frac{3\pi}{4} = \frac{3\pi}{4} - \frac{3\pi}{4} = 0$.
$x_2 = 10\cos(0) = 10\,\mathrm{cm}$.
$v_2 = -40\pi\sin(0) = 0\,\mathrm{cm/s}$.
Quãng đường vật đi được từ $t_1$ đến $t_2$ là $\Delta x = |x_2 - x_1| = |10 - 0| = 10\,\mathrm{cm}$.
Thời gian là $\Delta t = t_2 - t_1 = 0,1875 - 0,0625 = 0,125\,\mathrm{s}$.
Tốc độ trung bình là $v_{tb} = \frac{\Delta x}{\Delta t} = \frac{10}{0,125} = 80\,\mathrm{cm/s}$
}
\end{ex}

% ===== Câu 99 =====
% ID: L11C1B1-127-TN
% Mức: TH
\begin{ex}
% ID: L11C1B1-127-TN
% Mức: TH
Một vật dao động điều hòa với phương trình
$x=0,05\cos\left(20t+\dfrac{\pi}{2}\right) (\mathrm{m})$
với $t$ đo bằng giây. Vận tốc trung bình trong $\dfrac{1}{4}$ chu kì kể từ lúc $t=0$ là
\choice
{$-\pi\,\mathrm{m/s}$}
{$\dfrac{2}{\pi}\,\mathrm{m/s}$}
{\True $-\dfrac{2}{\pi}\,\mathrm{m/s}$}
{$\pi\,\mathrm{m/s}$}
\loigiai{
Vận tốc trung bình trong một khoảng thời gian được tính bằng tỉ số giữa độ dời và khoảng thời gian đó.
Phương trình dao động là $x=0,05\cos\left(20t+\dfrac{\pi}{2}\right) (\mathrm{m})$.
Tại thời điểm $t=0$, vị trí của vật là $x_0 = 0,05\cos\left(20(0)+\dfrac{\pi}{2}\right) = 0,05\cos\left(\dfrac{\pi}{2}\right) = 0 \, \mathrm{m}$.

hu kì dao động là $T = \dfrac{2\pi}{\omega} = \dfrac{2\pi}{20} = \dfrac{\pi}{10} \, \mathrm{s}$.
Khoảng thời gian xét là $\Delta t = \dfrac{1}{4}T = \dfrac{1}{4} \times \dfrac{\pi}{10} = \dfrac{\pi}{40} \, \mathrm{s}$.
Thời điểm kết thúc khoảng thời gian này là $t_1 = t_0 + \Delta t = 0 + \dfrac{\pi}{40} = \dfrac{\pi}{40} \, \mathrm{s}$.
Tại thời điểm $t_1 = \dfrac{\pi}{40} \, \mathrm{s}$, vị trí của vật là $x_1 = 0,05\cos\left(20\left(\dfrac{\pi}{40}\right)+\dfrac{\pi}{2}\right) = 0,05\cos\left(\dfrac{\pi}{2}+\dfrac{\pi}{2}\right) = 0,05\cos(\pi) = -0,05 \, \mathrm{m}$.
Độ dời của vật trong khoảng thời gian $\Delta t$ là $\Delta x = x_1 - x_0 = -0,05 - 0 = -0,05 \, \mathrm{m}$.
Vận tốc trung bình trong $\dfrac{1}{4}$ chu kì kể từ lúc $t=0$ là $v_{tb} = \dfrac{\Delta x}{\Delta t} = \dfrac{-0,05}{\dfrac{\pi}{40}} = \dfrac{-0,05 \times 40}{\pi} = \dfrac{-2}{\pi} \, \mathrm{m/s}$.
Vận tốc trung bình trong $\dfrac{1}{4}$ chu kì kể từ lúc $t=0$ là $-\dfrac{2}{\pi}\,\mathrm{m/s}$
}
\end{ex}

% ===== Câu 100 =====
% ID: L11C1B1-128-TN
% Mức: TH
\begin{ex}
% ID: L11C1B1-128-TN
% Mức: TH
Một chất điểm dao động điều hòa dạng hàm cos có chu kì $T$, biên độ $A$. Tốc độ trung bình của chất điểm khi pha của dao động biến thiên từ $-\dfrac{\pi}{2}$ đến $0$ bằng
\choice
{$\dfrac{3A}{T}$}
{\True $\dfrac{4A}{T}$}
{$\dfrac{3,6A}{T}$}
{$\dfrac{2A}{T}$}
\loigiai{
Đáp án B.
}
\end{ex}

% ===== Câu 101 =====
% ID: L11C1B1-129-TN
% Mức: TH
\begin{ex}
% ID: L11C1B1-129-TN
% Mức: TH
Một vật nhỏ dao động điều hòa dọc theo trục $Ox$, với $O$ là vị trí cân bằng, có phương trình dao động
$x=2\cos\left(2\pi t-\dfrac{\pi}{12}\right) (\mathrm{cm})$
với $t$ tính bằng giây. Tốc độ trung bình của vật từ thời điểm $t_1=\dfrac{13}{6}\,\mathrm{s}$ đến thời điểm $t_2=\dfrac{11}{3}\,\mathrm{s}$ gần nhất với giá trị nào sau đây?
\choice
{$11\,\mathrm{cm/s}$}
{$12\,\mathrm{cm/s}$}
{$54\,\mathrm{cm/s}$}
{\True $7\,\mathrm{cm/s}$}
\loigiai{
Đáp án D.
}
\end{ex}

% ===== Câu 102 =====
% ID: L11C1B1-130-TN
% Mức: TH
\begin{ex}
% ID: L11C1B1-130-TN
% Mức: TH
Một con lắc lò xo dao động với phương trình
$x=4\cos\left(4\pi t-\dfrac{\pi}{8}\right) (\mathrm{cm})$
với $t$ đo bằng giây. Tốc độ trung bình của vật từ thời điểm $t_1=0,03125\,\mathrm{s}$ đến thời điểm $t_2=2,90625\,\mathrm{s}$ gần nhất với giá trị nào sau đây?
\choice
{$11\,\mathrm{cm/s}$}
{$12\,\mathrm{cm/s}$}
{$54\,\mathrm{cm/s}$}
{\True $27\,\mathrm{cm/s}$}
\loigiai{
Đáp án D.
}
\end{ex}

% ===== Câu 103 =====
% ID: L11C1B1-131-TN
% Mức: TH
\begin{ex}
% ID: L11C1B1-131-TN
% Mức: TH
Một chất điểm dao động điều hòa với phương trình
$x=6\cos(10\pi t) (\mathrm{cm})$
với $t$ đo bằng giây. Vận tốc trung bình của chất điểm sau $\dfrac{1}{4}$ chu kì tính từ $t=0$ là
\choice
{$+1,2\,\mathrm{m/s}$}
{\True $-1,2\,\mathrm{m/s}$}
{$-2\,\mathrm{m/s}$}
{$+2\,\mathrm{m/s}$}
\loigiai{
Đáp án B.
}
\end{ex}

% ===== Câu 104 =====
% ID: L11C1B1-132-TN
% Mức: TH
\begin{ex}
% ID: L11C1B1-132-TN
% Mức: TH
Một vật dao động điều hòa theo phương trình
$x=5\cos\left(4\pi t+\dfrac{\pi}{3}\right) (\mathrm{cm})$
với $t$ đo bằng giây. Tốc độ trung bình và vận tốc trung bình của vật trong khoảng thời gian tính từ lúc $t=0$ đến thời điểm vật đi qua vị trí cân bằng theo chiều dương lần thứ nhất lần lượt là
\choice
{$60\,\mathrm{cm/s}$ và $8,6\,\mathrm{cm/s}$}
{\True $42,9\,\mathrm{cm/s}$ và $-8,6\,\mathrm{cm/s}$}
{$42,9\,\mathrm{cm/s}$ và $8,6\,\mathrm{cm/s}$}
{$30\,\mathrm{cm/s}$ và $8,6\,\mathrm{cm/s}$}
\loigiai{
Đáp án B.
}
\end{ex}

% ===== Câu 105 =====
% ID: L11C1B1-133-TN
% Mức: TH
\begin{ex}
% ID: L11C1B1-133-TN
% Mức: TH
Một chất điểm dao động điều hòa theo phương trình
$x=8\cos\left(2\pi t-\dfrac{\pi}{6}\right) (\mathrm{cm})$
với $t$ đo bằng giây. Tốc độ trung bình nhỏ nhất mà chất điểm đạt được trong khoảng thời gian $\dfrac{4}{3}\,\mathrm{s}$ là
\choice
{\True $30\,\mathrm{cm/s}$}
{$36\,\mathrm{cm/s}$}
{$24\,\mathrm{cm/s}$}
{$6\,\mathrm{cm/s}$}
\loigiai{
Đáp án A.
}
\end{ex}

% ===== Câu 106 =====
% ID: L11C1B1-134-TN
% Mức: TH
\begin{ex}
% ID: L11C1B1-134-TN
% Mức: TH
Một vật dao động điều hòa dọc theo trục $Ox$ với phương trình
$x=3\cos\left(4\pi t-\dfrac{\pi}{3}\right) (\mathrm{cm})$
với $t$ đo bằng giây. Tốc độ trung bình của vật từ thời điểm $t_1=\dfrac{13}{6}\,\mathrm{s}$ đến thời điểm $t_2=\dfrac{23}{6}\,\mathrm{s}$ là
\choice
{$16,2\,\mathrm{cm/s}$}
{$40,54\,\mathrm{cm/s}$}
{\True $24,3\,\mathrm{cm/s}$}
{$45\,\mathrm{cm/s}$}
\loigiai{
Đáp án C.
}
\end{ex}

% ===== Câu 107 =====
% ID: L11C1B1-135-TN
% Mức: TH
\begin{ex}
% ID: L11C1B1-135-TN
% Mức: TH
Một vật dao động điều hòa với biên độ $6\,\mathrm{cm}$. Thời gian ngắn nhất mà vật đi từ vị trí cân bằng đến vị trí có động năng bằng $3$ lần thế năng là $0,1\,\mathrm{s}$. Tốc độ trung bình của con lắc trong nửa chu kì là
\choice
{$5\,\mathrm{cm/s}$}
{$10\,\mathrm{cm/s}$}
{\True $20\,\mathrm{cm/s}$}
{$15\,\mathrm{cm/s}$}
\loigiai{
Đáp án C.
}
\end{ex}

% ===== Câu 108 =====
% ID: L11C1B1-136-TN
% Mức: TH
\begin{ex}
% ID: L11C1B1-136-TN
% Mức: TH
Một chất điểm dao động điều hòa dọc theo trục $Ox$ với phương trình
$x=8\cos\left(4\pi t+\dfrac{\pi}{6}\right) (\mathrm{cm})$
với $t$ đo bằng giây. Tốc độ trung bình của vật từ thời điểm $t_1=2,375\,\mathrm{s}$ đến thời điểm $t_2=4,75\,\mathrm{s}$ gần nhất với giá trị nào sau đây?
\choice
{\True $49\,\mathrm{cm/s}$}
{$40,54\,\mathrm{cm/s}$}
{$54,9\,\mathrm{cm/s}$}
{$45\,\mathrm{cm/s}$}
\loigiai{
Đáp án A.
}
\end{ex}

% ===== Câu 109 =====
% ID: L11C1B1-137-TN
% Mức: TH
\begin{ex}
% ID: L11C1B1-137-TN
% Mức: TH
Một chất điểm dao động điều hòa theo phương trình
$x=5\cos(20\pi t) (\mathrm{cm})$
với $t$ đo bằng giây. Tốc độ trung bình lớn nhất mà chất điểm đạt được trong khoảng thời gian $\dfrac{1}{6}$ chu kì là
\choice
{$100\,\mathrm{cm/s}$}
{$50\pi\,\mathrm{cm/s}$}
{$100\pi\,\mathrm{cm/s}$}
{\True $300\,\mathrm{cm/s}$}
\loigiai{
Đáp án D.
}
\end{ex}

% ===== Câu 110 =====
% ID: L11C1B1-138-TN
% Mức: TH
\begin{ex}
% ID: L11C1B1-138-TN
% Mức: TH
Vật dao động điều hòa với tần số $f=0,5\,\mathrm{Hz}$. Tại $t=0$, vật có li độ $x=4\,\mathrm{cm}$ và vận tốc $v=-4\pi\,\mathrm{cm/s}$. Tốc độ trung bình của vật từ thời điểm $t_1=0$ đến thời điểm $t_2=2,5\,\mathrm{s}$ gần nhất với giá trị nào sau đây?
\choice
{$11\,\mathrm{cm/s}$}
{\True $12\,\mathrm{cm/s}$}
{$54\,\mathrm{cm/s}$}
{$15\,\mathrm{cm/s}$}
\loigiai{
Đáp án B.
}
\end{ex}

% ===== Câu 111 =====
% ID: L11C1B1-139-TN
% Mức: TH
\begin{ex}
% ID: L11C1B1-139-TN
% Mức: TH
Một con lắc lò xo gồm lò xo có khối lượng không đáng kể, độ cứng $50\,\mathrm{N/m}$ và vật $M$ có khối lượng $200\,\mathrm{g}$ trượt không ma sát trên mặt phẳng nằm ngang. Kéo $M$ ra khỏi vị trí cân bằng một đoạn $4\,\mathrm{cm}$ rồi buông nhẹ thì vật dao động điều hòa. Tốc độ trung bình của $M$ sau khi nó đi được quãng đường $6\,\mathrm{cm}$ kể từ khi bắt đầu chuyển động là bao nhiêu? Lấy $\pi^2=10$.
\choice
{$60\,\mathrm{cm/s}$}
{\True $45\,\mathrm{cm/s}$}
{$40\,\mathrm{cm/s}$}
{$30\,\mathrm{cm/s}$}
\loigiai{
Đáp án B.
}
\end{ex}

% ===== Câu 112 =====
% ID: L11C1B1-140-TN
% Mức: TH
\begin{ex}
% ID: L11C1B1-140-TN
% Mức: TH
Một chất điểm dao động điều hòa với chu kì $T$. Trong khoảng thời gian ngắn nhất khi đi từ vị trí có li độ $x=\dfrac{A}{2}$ đến vị trí có li độ $x=-\dfrac{A}{2}$, chất điểm có tốc độ trung bình là
\choice
{\True $\dfrac{6A}{T}$}
{$\dfrac{4,5A}{T}$}
{$\dfrac{1,5A}{T}$}
{$\dfrac{4A}{T}$}
\loigiai{
Đáp án A.
}
\end{ex}

% ===== Câu 113 =====
% ID: L11C1B1-141-TN
% Mức: TH
\begin{ex}
% ID: L11C1B1-141-TN
% Mức: TH
Một vật dao động điều hòa với biên độ $A$ và chu kì $T=0,4\,\mathrm{s}$. Khi vật có li độ $1,2\,\mathrm{cm}$ thì động năng chiếm $96\%$ cơ năng. Tốc độ trung bình trong một chu kì là
\choice
{$1,2\,\mathrm{m/s}$}
{$0,3\,\mathrm{m/s}$}
{$0,2\,\mathrm{m/s}$}
{\True $0,6\,\mathrm{m/s}$}
\loigiai{
Đáp án D.
}
\end{ex}

% ===== Câu 114 =====
% ID: L11C1B1-142-TN
% Mức: TH
\begin{ex}
% ID: L11C1B1-142-TN
% Mức: TH
Một vật dao động điều hòa với biên độ $A=10\,\mathrm{cm}$ và chu kì $T=0,2\,\mathrm{s}$. Tốc độ trung bình lớn nhất của vật trong khoảng thời gian $\Delta t=\dfrac{1}{15}\,\mathrm{s}$ là
\choice
{$1,5\,\mathrm{m/s}$}
{$1,3\,\mathrm{m/s}$}
{$2,1\,\mathrm{m/s}$}
{\True $2,6\,\mathrm{m/s}$}
\loigiai{
Đáp án D.
}
\end{ex}

% ===== Câu 115 =====
% ID: L11C1B1-143-TN
% Mức: TH
\begin{ex}
% ID: L11C1B1-143-TN
% Mức: TH
Một vật dao động với chu kì $4\,\mathrm{s}$ trên quỹ đạo có chiều dài $2\,\mathrm{cm}$ theo phương trình
$x=A\cos\left(\omega t+\dfrac{\pi}{4}\right) (\mathrm{cm}).$
Vận tốc trung bình của vật sau $3\,\mathrm{s}$ là
\choice
{$0,5\,\mathrm{cm/s}$}
{$-1\,\mathrm{cm/s}$}
{\True $0\,\mathrm{cm/s}$}
{$-1,4\,\mathrm{cm/s}$}
\loigiai{
Đáp án C.
}
\end{ex}

% ===== Câu 116 =====
% ID: L11C1B1-144-TN
% Mức: TH
\begin{ex}
% ID: L11C1B1-144-TN
% Mức: TH
Một vật dao động điều hòa với chu kì $T$ và biên độ $A$. Tốc độ trung bình lớn nhất của vật thực hiện được trong khoảng thời gian $\dfrac{2T}{3}$ là
\choice
{\True $\dfrac{4,5A}{T}$}
{$\dfrac{6A}{T}$}
{$\dfrac{\sqrt{3}A}{T}$}
{$\dfrac{1,5\sqrt{3}A}{T}$}
\loigiai{
Đáp án A.
}
\end{ex}

% ===== Câu 117 =====
% ID: L11C1B1-145-TN
% Mức: TH
\begin{ex}
% ID: L11C1B1-145-TN
% Mức: TH
Một chất điểm dao động điều hòa trên đoạn đường $PQ=20\,\mathrm{cm}$, thời gian vật đi từ $P$ đến $Q$ là $0,5\,\mathrm{s}$. Gọi $O,E,F$ lần lượt là trung điểm của $PQ,OP$ và $OQ$. Tốc độ trung bình của chất điểm trên đoạn $EF$ là
\choice
{$1,2\,\mathrm{m/s}$}
{$0,8\,\mathrm{m/s}$}
{\True $0,6\,\mathrm{m/s}$}
{$0,4\,\mathrm{m/s}$}
\loigiai{
Đáp án C.
}
\end{ex}

% ===== Câu 118 =====
% ID: L11C1B1-146-TN
% Mức: NB
\dangbt{Xác định và tính toán các đại lượng vận tốc, gia tốc tại thời điểm xác định}
\begin{ex}
% ID: L11C1B1-146-TN
% Mức: NB
Vật đang dao động điều hòa dọc theo đường thẳng. Một điểm $M$ nằm cố định trên đường thẳng đó, phía ngoài khoảng chuyển động của vật. Tại thời điểm $t$ thì vật xa điểm $M$ nhất, sau đó một khoảng thời gian ngắn nhất là $\Delta t$ thì vật gần điểm $M$ nhất. Độ lớn vận tốc của vật sẽ đạt cực đại vào thời điểm gần nhất là
\choice
{$t+\Delta t$}
{\True $t+0,5\Delta t$}
{$0,5(t+\Delta t)$}
{$0,5t+0,25\Delta t$}
\loigiai{
Khi vật xa $M$ nhất và sau đó gần $M$ nhất thì vật đi từ biên này sang biên kia.
Do đó
$\Delta t=\dfrac{T}{2}.$
Tốc độ cực đại khi vật qua vị trí cân bằng, tức là ở giữa hai biên.
Vì vậy thời điểm gần nhất là
$t+\dfrac{\Delta t}{2}=t+0,5\Delta t.$
}
\end{ex}

% ===== Câu 119 =====
% ID: L11C1B1-147-TN
% Mức: NB
\begin{ex}
% ID: L11C1B1-147-TN
% Mức: NB
Vật đang dao động điều hòa với biên độ $A$ dọc theo đường thẳng. Một điểm $M$ nằm cố định trên đường thẳng đó, phía ngoài khoảng chuyển động của vật. Tại thời điểm $t$ thì vật gần điểm $M$ nhất, sau đó một khoảng thời gian ngắn nhất là $\Delta t$ thì vật xa điểm $M$ nhất. Vật cách vị trí cân bằng một khoảng $\dfrac{A}{\sqrt{2}}$ vào thời điểm gần nhất là
\choice
{$t+\dfrac{\Delta t}{3}$}
{$t+\dfrac{\Delta t}{6}$}
{\True $t+\dfrac{\Delta t}{4}$}
{$0,5t+0,25\Delta t$}
\loigiai{
Từ lúc vật gần $M$ nhất đến lúc xa $M$ nhất là từ một biên đến biên đối diện:
$\Delta t=\dfrac{T}{2}.$
Từ biên đến vị trí có
$|x|=\dfrac{A}{\sqrt{2}}$
ứng với góc quét
$\Delta\varphi=\dfrac{\pi}{4}.$
Thời gian tương ứng:
$t_1=\dfrac{\pi/4}{2\pi}T=\dfrac{T}{8}.$
Vì
$\Delta t=\dfrac{T}{2}$
nên
$t_1=\dfrac{\Delta t}{4}.$
Vậy thời điểm gần nhất là
$t+\dfrac{\Delta t}{4}.$
}
\end{ex}

% ===== Câu 120 =====
% ID: L11C1B1-148-TN
% Mức: NB
\begin{ex}
% ID: L11C1B1-148-TN
% Mức: NB
Một dao động điều hòa có chu kì dao động là $T$ và biên độ là $A$. Tại thời điểm ban đầu vật có li độ $x_0\gt 0$. Thời gian ngắn nhất để vật đi từ vị trí ban đầu về vị trí cân bằng gấp bốn lần thời gian ngắn nhất để vật đi từ vị trí ban đầu về vị trí biên $x=+A$. Chọn phương án đúng.
\choice
{$x_0=0{,}924A$}
{$x_0=0,5A\sqrt{3}$}
{\True $x_0=0{,}95A$}
{$x_0=0{,}022A$}
\loigiai{
Đặt
$x_0=A\cos\alpha,\qquad 0\lt \alpha<\dfrac{\pi}{2}.$
Thời gian ngắn nhất đi từ $x_0$ đến biên dương $x=+A$ là
$t_1=\dfrac{\alpha}{\omega}.$
Thời gian ngắn nhất đi từ $x_0$ đến vị trí cân bằng là
$t_2=\dfrac{\frac{\pi}{2}-\alpha}{\omega}.$
Theo đề:
$t_2=4t_1.$
Suy ra
 $\dfrac{\pi}{2}-\alpha=4\alpha\Rightarrow\alpha=\dfrac{\pi}{10}$. 
Do đó
 x_0= $A\cos\dfrac{\pi}{10}\approx0,95\text{A}$.
}
\end{ex}

% ===== Câu 121 =====
% ID: L11C1B1-149-TN
% Mức: NB
\begin{ex}
% ID: L11C1B1-149-TN
% Mức: NB
Một vật dao động điều hòa với chu kì $T$ trên đoạn thẳng $PQ$. Gọi $O$, $E$ lần lượt là trung điểm của $PQ$ và $OQ$. Thời gian để vật đi từ $O$ đến $Q$ rồi đến $E$ là
\choice
{$\dfrac{5T}{6}$}
{\True $\dfrac{5T}{12}$}
{$\dfrac{T}{12}$}
{$\dfrac{7T}{12}$}
\loigiai{
Vì $O$ là vị trí cân bằng, $Q$ là biên dương nên thời gian đi từ $O$ đến $Q$ là $t_{OQ}=\dfrac{T}{4}$.

Vì $E$ là trung điểm của $OQ$ nên $OE=\dfrac{A}{2}$.

Thời gian đi từ $Q$ về $E$ tương ứng từ $x=A$ đến $x=\dfrac{A}{2}$: $t_{QE}=\dfrac{T}{6}$.

Vậy $t=t_{OQ}+t_{QE}=\dfrac{T}{4}+\dfrac{T}{6}=\dfrac{5T}{12}$.
}
\end{ex}

% ===== Câu 122 =====
% ID: L11C1B1-150-TN
% Mức: NB
\begin{ex}
% ID: L11C1B1-150-TN
% Mức: NB
Một con lắc lò xo có vật nặng với khối lượng $m=100\ \text{g}$ và lò xo có độ cứng $k=10\ \text{N/m}$ đang dao động điều hòa với biên độ $2\ \text{cm}$. Trong mỗi chu kì dao động, thời gian mà vật nặng ở cách vị trí cân bằng lớn hơn $1\ \text{cm}$ là
\choice
{$0,32\ \text{s}$}
{$0,22\ \text{s}$}
{\True $0,42\ \text{s}$}
{$0,52\ \text{s}$}
\loigiai{
Ta có $m = 100\ \mathrm{g} = 0,1\ \mathrm{kg}$.

Tần số góc:
$\omega = \sqrt{\dfrac{k}{m}} = \sqrt{\dfrac{10}{0,1}} = 10\ \mathrm{rad/s}$.

Chu kì:
$T = \dfrac{2\pi}{\omega} = \dfrac{2\pi}{10} \approx 0,628\ \mathrm{s}$.

Điều kiện vật cách vị trí cân bằng lớn hơn $1\ \mathrm{cm}$ là $|x| \gt  1 = \dfrac{A}{2}$.

Trong một chu kì, thời gian thỏa điều kiện $|x| \gt  \dfrac{A}{2}$ là $\Delta t = \dfrac{2T}{3}$.

Suy ra $\Delta t = \dfrac{2}{3} \cdot 0,628 \approx 0,42\ \mathrm{s}$.
}
\end{ex}

% ===== Câu 123 =====
% ID: L11C1B1-151-TN
% Mức: NB
\begin{ex}
% ID: L11C1B1-151-TN
% Mức: NB
Một vật dao động điều hòa có phương trình li độ
$x=8\cos\left(7\pi t+\dfrac{\pi}{6}\right)\ \text{cm}.$
Khoảng thời gian tối thiểu để vật đi từ li độ $4\ \text{cm}$ đến vị trí có li độ $-4\sqrt{3}\ \text{cm}$ là
\choice
{$\dfrac{1}{24}\ \text{s}$}
{$\dfrac{5}{12}\ \text{s}$}
{\True $\dfrac{1}{14}\ \text{s}$}
{$\dfrac{1}{12}\ \text{s}$}
\loigiai{
Biên độ dao động là
$A=8\ \text{cm},\qquad \omega=7\pi\ \text{rad/s}.$
Ta có
$x_1=4=\dfrac{A}{2},\qquad x_2=-4\sqrt{3}=-\dfrac{A\sqrt{3}}{2}.$
Trên vòng tròn lượng giác, góc quét ngắn nhất từ $x=\dfrac{A}{2}$ đến $x=-\dfrac{A\sqrt{3}}{2}$ là
$\Delta\varphi=\dfrac{\pi}{2}.$
Do đó
$\Delta t= \dfrac{\Delta\varphi}{\omega} = \dfrac{\pi/2}{7\pi} = \dfrac{1}{14}\ \text{s}.$
}
\end{ex}

% ===== Câu 124 =====
% ID: L11C1B1-152-TN
% Mức: NB
\begin{ex}
% ID: L11C1B1-152-TN
% Mức: NB
Một vật dao động điều hòa với phương trình
$x=8\cos 2\pi t\ \text{cm},$
trong đó $t$ đo bằng giây. Vật phải mất thời gian tối thiểu bao nhiêu giây để đi từ vị trí $x=+8\ \text{cm}$ về vị trí $x=4\ \text{cm}$ mà vectơ vận tốc cùng hướng với hướng của trục tọa độ?
\choice
{$\dfrac{1}{3}\ \text{s}$}
{\True $\dfrac{5}{6}\ \text{s}$}
{$\dfrac{1}{2}\ \text{s}$}
{$\dfrac{1}{6}\ \text{s}$}
\loigiai{
Ta có
$x=8\cos 2\pi t,\qquad A=8\ \text{cm}.$
Vị trí $x=4\ \text{cm}$ ứng với
$\cos 2\pi t=\dfrac{1}{2}.$
Các pha thỏa mãn là
$2\pi t=\dfrac{\pi}{3}\quad \text{hoặc}\quad 2\pi t=\dfrac{5\pi}{3}.$
Vận tốc:
$v=x'=-16\pi \sin 2\pi t.$
Để vận tốc cùng hướng dương của trục tọa độ thì $v>0$, suy ra
$\sin 2\pi t<0.$
Vậy chọn
 2 $\pi$ t= $\dfrac{5\pi}{3}\Rightarrow$ t= $\dfrac{5}{6}\\text{s}$.
}
\end{ex}

% ===== Câu 125 =====
% ID: L11C1B1-153-TN
% Mức: TH
\begin{ex}
% ID: L11C1B1-153-TN
% Mức: TH
Vật đang dao động điều hòa với biên độ $A$ dọc theo đường thẳng. Một điểm $M$ nằm cố định trên đường thẳng đó, phía ngoài khoảng chuyển động của vật. Tại thời điểm $t$ thì vật xa điểm $M$ nhất, sau đó một khoảng thời gian ngắn nhất là $\Delta t$ thì vật gần điểm $M$ nhất. Vật cách vị trí cân bằng một khoảng $0,5A$ vào thời điểm gần nhất là
\choice
{\True $t+\dfrac{\Delta t}{3}$}
{$t+\dfrac{\Delta t}{6}$}
{$0,5(t+\Delta t)$}
{$0,5t+0,25\Delta t$}
\loigiai{
Từ lúc vật xa $M$ nhất đến lúc gần $M$ nhất là từ một biên đến biên đối diện, nên
$\Delta t=\dfrac{T}{2}.$
Từ biên đến vị trí có
$|x|=\dfrac{A}{2}$
góc quét là
$\Delta\varphi=\dfrac{\pi}{3}.$
Thời gian tương ứng:
$t_1=\dfrac{\pi/3}{2\pi}T=\dfrac{T}{6}.$
Mà $\Delta t=\dfrac{T}{2}$ nên
$t_1=\dfrac{\Delta t}{3}.$
Vậy thời điểm gần nhất là
$t+\dfrac{\Delta t}{3}.$
}
\end{ex}

% ===== Câu 126 =====
% ID: L11C1B1-154-TN
% Mức: VD
\begin{ex}
% ID: L11C1B1-154-TN
% Mức: VD
Một chất điểm dao động điều hòa theo quỹ đạo thẳng có chiều dài $8\ \text{cm}$. Thời gian ngắn nhất để vật đi từ vị trí có li độ $x_1=4\ \text{cm}$ đến $x_2=-2\sqrt{3}\ \text{cm}$ là $2\ \text{s}$. Tốc độ cực đại của vật trong quá trình dao động là
\choice
{$\dfrac{T}{8}$}
{$\dfrac{T}{16}$}
{\True $\dfrac{T}{6}$}
{$\dfrac{T}{12}$}
\loigiai{
Câu này đang bị lệch giữa đề và phương án.

Nếu hỏi tốc độ cực đại thì các phương án phải có đơn vị tốc độ, ví dụ $\text{cm/s}$, trong khi các phương án đang là các khoảng thời gian theo $T$.

Theo dữ kiện:
$2A = 8\ \text{cm} \Rightarrow A = 4\ \text{cm}.$

Vị trí $x_1 = 4 = A$, $x_2 = -2\sqrt{3} = -\dfrac{A\sqrt{3}}{2}.$

Góc quét ngắn nhất từ $x = A$ đến $x = -\dfrac{A\sqrt{3}}{2}$ là $\Delta\varphi = \dfrac{5\pi}{6}.$

Vì thời gian tương ứng là $2\ \text{s}$ nên $\omega = \dfrac{\Delta\varphi}{\Delta t} = \dfrac{5\pi/6}{2} = \dfrac{5\pi}{12}\ \mathrm{rad/s}.$

Tốc độ cực đại: $v_{\max} = \omega A = \dfrac{5\pi}{12} \cdot 4 = \dfrac{5\pi}{3}\ \text{cm/s}.$

Vậy nếu đúng là hỏi tốc độ cực đại thì đáp án phải là $v_{\max} = \dfrac{5\pi}{3}\ \text{cm/s}.$

Các phương án hiện tại không phù hợp với câu hỏi. Dấu True ở phương án $C$ chỉ được giữ theo bảng đáp án thầy cung cấp, không khớp với nội dung đề hiện tại.
}
\end{ex}

% ===== Câu 127 =====
% ID: L11C1B1-155-TN
% Mức: VD
\begin{ex}
% ID: L11C1B1-155-TN
% Mức: VD
Một chất điểm dao động điều hòa với chu kì $T$ trên trục $Ox$, với $O$ là vị trí cân bằng. Thời gian ngắn nhất vật đi từ li độ $x=\dfrac{A}{\sqrt{2}}$ đến li độ $x=\dfrac{A}{2}$ là
\choice
{\True $\dfrac{T}{24}$}
{$\dfrac{T}{16}$}
{$\dfrac{T}{6}$}
{$\dfrac{T}{12}$}
\loigiai{
Ta có
$\dfrac{A}{\sqrt{2}}=\dfrac{A\sqrt{2}}{2},\qquad \dfrac{A}{2}.$
Góc ứng với $x=\dfrac{A\sqrt{2}}{2}$ là $\dfrac{\pi}{4}$, góc ứng với $x=\dfrac{A}{2}$ là $\dfrac{\pi}{3}$.
Do đó góc quét ngắn nhất:
$\Delta\varphi=\dfrac{\pi}{3}-\dfrac{\pi}{4}=\dfrac{\pi}{12}.$
Suy ra
 $\Delta$ t= $\dfrac{\Delta\varphi}{2\pi}T$ = $\dfrac{\pi/12}{2\pi}T$ = $\dfrac{T}{24}$.
}
\end{ex}

% ===== Câu 128 =====
% ID: L11C1B1-156-TN
% Mức: VD
\begin{ex}
% ID: L11C1B1-156-TN
% Mức: VD
Một chất điểm đang dao động điều hòa trên một đoạn thẳng. Trên đoạn thẳng đó có năm điểm theo đúng thứ tự $M$, $N$, $O$, $P$ và $Q$, với $O$ là vị trí cân bằng. Biết cứ $0,05\ \text{s}$ thì chất điểm lại đi qua các điểm $M$, $N$, $O$, $P$ và $Q$; tốc độ tại $M$ và $Q$ bằng $0$. Chu kì bằng
\choice
{$0,3\ \text{s}$}
{\True $0,4\ \text{s}$}
{$0,2\ \text{s}$}
{$0,1\ \text{s}$}
\loigiai{
Vì tốc độ tại $M$ và $Q$ bằng $0$ nên $M$, $Q$ là hai vị trí biên.
Từ $M$ đến $Q$ là nửa chu kì, gồm $4$ khoảng thời gian bằng nhau:
$\dfrac{T}{2}=4\cdot 0,05=0,20\ \text{s}.$
Suy ra
$T=0,40\ \text{s}.$
}
\end{ex}

% ===== Câu 129 =====
% ID: L11C1B1-157-TN
% Mức: VD
\begin{ex}
% ID: L11C1B1-157-TN
% Mức: VD
Một chất điểm dao động điều hòa. Khoảng thời gian trong một chu kì để vật cách vị trí cân bằng một khoảng nhỏ hơn nửa biên độ là $1\ \text{s}$. Chu kì dao động là
\choice
{\True $3\ \text{s}$}
{$1,5\ \text{s}$}
{$6\ \text{s}$}
{$2\ \text{s}$}
\loigiai{
Từ kết quả quen thuộc:
$|x|\lt \dfrac{A}{2}$
trong một chu kì ứng với thời gian
$\Delta t=\dfrac{T}{3}.$
Theo đề:
 $\dfrac{T}{3}=1\Rightarrow$ T=3 $\\text{s}$.
}
\end{ex}

% ===== Câu 130 =====
% ID: L11C1B1-158-TN
% Mức: VD
\begin{ex}
% ID: L11C1B1-158-TN
% Mức: VD
Một chất điểm dao động điều hòa với biên độ $4\ \text{cm}$ và chu kì $0,9\ \text{s}$. Khoảng thời gian ngắn nhất để nó đi từ vị trí có li độ $+3\ \text{cm}$ đến vị trí cân bằng là
\choice
{$0,1035\ \text{s}$}
{\True $0,1215\ \text{s}$}
{$6,9601\ \text{s}$}
{$5,9315\ \text{s}$}
\loigiai{
Ta có
$\frac{x}{A}=\frac{3}{4}=0,75.$
Từ $x=3\ \text{cm}$ đến vị trí cân bằng $x=0$, góc quét ngắn nhất là
$\Delta\varphi=\frac{\pi}{2}-\arccos 0,75.$
Suy ra
$\Delta t=\frac{\Delta\varphi}{2\pi}T=\frac{\frac{\pi}{2}-\arccos 0,75}{2\pi}\cdot 0,9\approx 0,1215\ \text{s}.$
}
\end{ex}

% ===== Câu 131 =====
% ID: L11C1B1-159-TN
% Mức: VD
\begin{ex}
% ID: L11C1B1-159-TN
% Mức: VD
Một chất điểm dao động điều hòa với biên độ $4\ \text{cm}$ và chu kì $0,9\ \text{s}$. Khoảng thời gian ngắn nhất để nó đi từ vị trí có li độ $+3\ \text{cm}$ đến li độ $+4\ \text{cm}$ là
\choice
{\True $0,1035\ \text{s}$}
{$0,1215\ \text{s}$}
{$6,9601\ \text{s}$}
{$5,9315\ \text{s}$}
\loigiai{
Ta có
$\dfrac{x}{A}=\dfrac{3}{4}.$
Từ $x=3\ \mathrm{cm}$ đến biên dương $x=A=4\ \mathrm{cm}$, góc quét ngắn nhất là
$\Delta\varphi=\arccos\dfrac{3}{4}.$
Suy ra
$\Delta t=\dfrac{\Delta\varphi}{2\pi}T=\dfrac{\arccos(3/4)}{2\pi}\cdot 0,9\approx 0,1035\ \mathrm{s}.$
}
\end{ex}

% ===== Câu 132 =====
% ID: L11C1B1-160-TN
% Mức: VD
\begin{ex}
% ID: L11C1B1-160-TN
% Mức: VD
Một chất điểm dao động điều hòa trên trục $Ox$ với chu kì $T$. Vị trí cân bằng của chất điểm trùng với gốc tọa độ. Khoảng thời gian ngắn nhất để nó đi từ vị trí có li độ $x=A$ đến vị trí có li độ $x=-\dfrac{A}{2}$ là
\choice
{$\dfrac{T}{8}$}
{$\dfrac{T}{6}$}
{$\dfrac{T}{4}$}
{\True $\dfrac{T}{3}$}
\loigiai{
Từ $x=A$ đến $x=-\dfrac{A}{2}$, trên vòng tròn lượng giác ta có
$\cos\Delta\varphi=-\dfrac{1}{2}.$
Góc quét ngắn nhất là
$\Delta\varphi=\dfrac{2\pi}{3}.$
Suy ra
 $\Delta$ t= $\dfrac{\Delta\varphi}{2\pi}T$ = $\dfrac{2\pi/3}{2\pi}T$ = $\dfrac{T}{3}$.
}
\end{ex}

% ===== Câu 133 =====
% ID: L11C1B1-161-TN
% Mức: VD
\begin{ex}
% ID: L11C1B1-161-TN
% Mức: VD
Một chất điểm dao động điều hòa trên đoạn đường $PQ$, thời gian vật đi từ $P$ đến $Q$ là $0,25\ \text{s}$. Gọi $O$, $E$ lần lượt là trung điểm của $PQ$ và $OQ$. Thời gian ngắn nhất vật đi từ $E$ đến $Q$ là
\choice
{$\dfrac{1}{24}\ \text{s}$}
{$\dfrac{1}{16}\ \text{s}$}
{$\dfrac{1}{6}\ \text{s}$}
{\True $\dfrac{1}{12}\ \text{s}$}
\loigiai{
Hai điểm $P, Q$ là hai vị trí biên nên thời gian đi từ $P$ đến $Q$ là nửa chu kì: $\dfrac{T}{2} = 0,25 \Rightarrow T = 0,5\ \mathrm{s}$.

Vì $E$ là trung điểm của $OQ$ nên $OE = \dfrac{A}{2}$.

Thời gian ngắn nhất đi từ $E$ đến $Q$ tương ứng đi từ $x = \dfrac{A}{2}$ đến $x = A$, có góc quét $\Delta\varphi = \dfrac{\pi}{3}$.

Vậy $\Delta t = \dfrac{\pi/3}{2\pi} \cdot 0,5 = \dfrac{1}{12}\ \mathrm{s}$.
}
\end{ex}

% ===== Câu 134 =====
% ID: L11C1B1-162-TN
% Mức: VD
\begin{ex}
% ID: L11C1B1-162-TN
% Mức: VD
Một chất điểm dao động điều hòa với chu kì $1\ \text{s}$, biên độ $4,5\ \text{cm}$. Khoảng thời gian trong một chu kì để vật cách vị trí cân bằng một khoảng lớn hơn $2\ \text{cm}$ là
\choice
{$0,29\ \text{s}$}
{$16,80\ \text{s}$}
{\True $0,71\ \text{s}$}
{$0,15\ \text{s}$}
\loigiai{
Điều kiện vật cách vị trí cân bằng lớn hơn $2\ \mathrm{cm}$ là $|x| \gt  2$.

Ta có $\dfrac{2}{A} = \dfrac{2}{4,5} \approx 0,444$.

Trong một chu kì, thời gian để $|x| \gt  2$ là $\Delta t = \dfrac{4\arccos(2/4,5)}{2\pi}T$.

Thay $T = 1\ \mathrm{s}$: $\Delta t = \dfrac{4\arccos(2/4,5)}{2\pi} \approx 0,71\ \mathrm{s}$.
}
\end{ex}

% ===== Câu 135 =====
% ID: L11C1B1-163-TN
% Mức: VD
\begin{ex}
% ID: L11C1B1-163-TN
% Mức: VD
Một vật dao động điều hòa với biên độ $A$, chu kì $T$. Thời gian ngắn nhất vật đi từ vị trí $x=-0{,}5A$ đến vị trí có $x=+0{,}5A$ là
\choice
{$\dfrac{T}{2}$}
{$\dfrac{T}{12}$}
{$\dfrac{T}{4}$}
{\True $\dfrac{T}{6}$}
\loigiai{
Từ $x=-\dfrac{A}{2}$ đến $x=\dfrac{A}{2}$ theo đường ngắn nhất qua vị trí cân bằng, góc quét là
$\Delta\varphi=\dfrac{\pi}{3}.$
Do đó
 $\Delta$ t= $\dfrac{\Delta\varphi}{2\pi}T$ = $\dfrac{\pi/3}{2\pi}T$ = $\dfrac{T}{6}$.
}
\end{ex}

% ===== Câu 136 =====
% ID: L11C1B1-164-TN
% Mức: VD
\begin{ex}
% ID: L11C1B1-164-TN
% Mức: VD
Một chất điểm dao động điều hòa với biên độ $10\ \text{cm}$ và tần số góc $10\ \text{rad/s}$. Khoảng thời gian ngắn nhất để nó đi từ vị trí có li độ $+3,5\ \text{cm}$ đến vị trí có li độ $+10\ \text{cm}$ là
\choice
{$0,036\ \text{s}$}
{\True $0,121\ \text{s}$}
{$2,049\ \text{s}$}
{$6,951\ \text{s}$}
\loigiai{
Ta có
$\dfrac{x}{A}=\dfrac{3,5}{10}=0,35.$
Gọi $\Delta\varphi$ là góc quét ngắn nhất từ vị trí $x=3,5\ \mathrm{cm}$ đến biên dương $x=A$. Khi đó
$\Delta\varphi=\arccos 0,35.$
Suy ra
$\Delta t= \dfrac{\Delta\varphi}{\omega} = \dfrac{\arccos 0,35}{10}\approx 0,121\ \mathrm{s}.$
}
\end{ex}

% ===== Câu 137 =====
% ID: L11C1B1-165-TN
% Mức: VD
\begin{ex}
% ID: L11C1B1-165-TN
% Mức: VD
Một chất điểm dao động điều hòa với chu kì $T$. Khoảng thời gian trong một chu kì để vật cách vị trí cân bằng một khoảng lớn hơn $0,5\sqrt{2}$ biên độ là
\choice
{$\dfrac{T}{3}$}
{$\dfrac{2T}{3}$}
{$\dfrac{T}{6}$}
{\True $\dfrac{T}{2}$}
\loigiai{
Điều kiện:
$|x|\gt \dfrac{A\sqrt{2}}{2}.$
Trong một chu kì, tổng góc quét thỏa điều kiện này là
$\Delta\varphi=\pi.$
Do đó
 $\Delta$ t= $\dfrac{\pi}{2\pi}T$ = $\dfrac{T}{2}$.
}
\end{ex}

% ===== Câu 13 =====
% ID: L11C1B1-166-DS
% Mức: TH
\dangbt{Viết phương trình dao động điều hòa.}
\begin{ex}
% ID: L11C1B1-166-DS
% Mức: TH
Vật có biên độ $4\,\text{cm}$, chu kì $1\,\text{s}$ và tại $t=0$ qua vị trí cân bằng theo chiều dương. Xét phương trình $x=4\cos(2\pi t/1-\pi/2)$ cm.
\choiceTF
{\True Biên độ là $4\,\text{cm}$.}
{\True Chu kì là $1\,\text{s}$.}
{\True Tại $t=0$, li độ bằng 0.}
{\True Tại $t=0$, vận tốc dương.}
\end{ex}

% ===== Câu 14 =====
% ID: L11C1B1-167-TL
% Mức: TH
\begin{bt}
% ID: L11C1B1-167-TL
% Mức: TH
Vật dao động điều hòa có biên độ $4\,\text{cm}$, chu kì $1\,\text{s}$; lúc $t=0$ đi qua vị trí cân bằng theo chiều âm. Viết phương trình li độ dạng cos.
\end{bt}

% ===== Câu 17 =====
% ID: L11C1B1-168-DS
% Mức: VD
\begin{ex}
% ID: L11C1B1-168-DS
% Mức: VD
Vật có biên độ $6\,\text{cm}$, chu kì $2\,\text{s}$ và tại $t=0$ qua vị trí cân bằng theo chiều dương. Xét phương trình $x=6\cos(2\pi t/2-\pi/2)$ cm.
\choiceTF
{\True Biên độ là $6\,\text{cm}$.}
{\True Chu kì là $2\,\text{s}$.}
{\True Tại $t=0$, li độ bằng 0.}
{\True Tại $t=0$, vận tốc dương.}
\end{ex}

% ===== Câu 18 =====
% ID: L11C1B1-169-TL
% Mức: VD
\begin{bt}
% ID: L11C1B1-169-TL
% Mức: VD
Vật dao động điều hòa có biên độ $6\,\text{cm}$, chu kì $2\,\text{s}$; lúc $t=0$ đi qua vị trí cân bằng theo chiều âm. Viết phương trình li độ dạng cos.
\end{bt}

% ===== Câu 21 =====
% ID: L11C1B1-170-DS
% Mức: VD
\begin{ex}
% ID: L11C1B1-170-DS
% Mức: VD
Vật có biên độ $8\,\text{cm}$, chu kì $4\,\text{s}$ và tại $t=0$ qua vị trí cân bằng theo chiều dương. Xét phương trình $x=8\cos(2\pi t/4-\pi/2)$ cm.
\choiceTF
{\True Biên độ là $8\,\text{cm}$.}
{\True Chu kì là $4\,\text{s}$.}
{\True Tại $t=0$, li độ bằng 0.}
{\True Tại $t=0$, vận tốc dương.}
\end{ex}

% ===== Câu 22 =====
% ID: L11C1B1-171-TL
% Mức: VD
\begin{bt}
% ID: L11C1B1-171-TL
% Mức: VD
Vật dao động điều hòa có biên độ $8\,\text{cm}$, chu kì $4\,\text{s}$; lúc $t=0$ đi qua vị trí cân bằng theo chiều âm. Viết phương trình li độ dạng cos.
\end{bt}

% ===== Câu 34 =====
% ID: L11C1B1-172-TLN
% Mức: TH
\begin{ex}
% ID: L11C1B1-172-TLN
% Mức: TH
Vật dao động điều hòa có biên độ $4\,\text{cm}$ và quỹ đạo dài bao nhiêu cm?
\shortans{7/1/1900}
\end{ex}

% ===== Câu 38 =====
% ID: L11C1B1-173-TLN
% Mức: VD
\begin{ex}
% ID: L11C1B1-173-TLN
% Mức: VD
Vật dao động điều hòa có biên độ $6\,\text{cm}$ và quỹ đạo dài bao nhiêu cm?
\shortans{12}
\end{ex}

% ===== Câu 42 =====
% ID: L11C1B1-174-TLN
% Mức: VD
\begin{ex}
% ID: L11C1B1-174-TLN
% Mức: VD
Vật dao động điều hòa có biên độ $8\,\text{cm}$ và quỹ đạo dài bao nhiêu cm?
\shortans{15/01/1900}
\end{ex}

% ===== Câu 53 =====
% ID: L11C1B1-175-TN
% Mức: TH
\begin{ex}
% ID: L11C1B1-175-TN
% Mức: TH
Vật có biên độ $9\,\text{cm}$, chu kì $1\,\text{s}$ và ở biên dương lúc $t=0$. Phương trình đúng là
\choice
{\True $x=9\cos(2\pi t/1)$ cm}
{$x=9\sin(2\pi t/1)$ cm}
{$x=9\cos(\pi t/1)$ cm}
{$x=9\cos(2\pi t/1+\pi/2)$ cm}
\end{ex}

% ===== Câu 55 =====
% ID: L11C1B1-176-TN
% Mức: TH
\begin{ex}
% ID: L11C1B1-176-TN
% Mức: TH
Một vật dao động điều hoà dọc theo trục $Ox$ với phương trình $x=A\cos \omega t$. Nếu chọn gốc toạ độ $O$ tại vị trí cân bằng của vật thì gốc thời gian $t=0$ là lúc vật
\choice
{\True ở vị trí li độ cực đại thuộc phần dương của trục Ox}
{qua vị trí cân bằng $O$ ngược chiều dương của trục Ox}
{ở vị trí li độ cực đại thuộc phần âm của trục Ox}
{qua vị trí cân bằng $O$ theo chiều dương của trục Ox}
\loigiai{
Gốc thời gian $t=0$ là lúc vật ở vị trí li độ cực đại thuộc phần dương của trục Ox
}
\end{ex}

% ===== Câu 57 =====
% ID: L11C1B1-177-TN
% Mức: VD
\begin{ex}
% ID: L11C1B1-177-TN
% Mức: VD
Vật có biên độ $10\,\text{cm}$, chu kì $2\,\text{s}$ và ở biên dương lúc $t=0$. Phương trình đúng là
\choice
{\True $x=10\cos(2\pi t/2)$ cm}
{$x=10\sin(2\pi t/2)$ cm}
{$x=10\cos(\pi t/2)$ cm}
{$x=10\cos(2\pi t/2+\pi/2)$ cm}
\end{ex}

